\documentclass[11pt]{amsart}
\usepackage{amssymb}
\usepackage{booktabs}
\usepackage[expansion=false]{microtype}
\usepackage{xcolor}
\usepackage{pgfplots}
\usepackage{xurl}
\pgfplotsset{compat=1.18}
\usepackage[colorlinks=true,linkcolor=blue!50!black,citecolor=blue!50!black,urlcolor=blue!50!black]{hyperref}

\newtheorem{theorem}{Theorem}[section]
\newtheorem{lemma}[theorem]{Lemma}
\newtheorem{proposition}[theorem]{Proposition}
\newtheorem{corollary}[theorem]{Corollary}
\theoremstyle{definition}
\newtheorem{definition}[theorem]{Definition}
\newtheorem{algorithm}[theorem]{Algorithm}
\theoremstyle{remark}
\newtheorem{remark}[theorem]{Remark}
\newtheorem{example}[theorem]{Example}

\newcommand{\Ein}{\operatorname{Ein}}
\newcommand{\Eip}{\mathcal{E}}
\newcommand{\erfc}{\operatorname{erfc}}
\newcommand{\erf}{\operatorname{erf}}
\newcommand{\erfi}{\operatorname{erfi}}
\newcommand{\Ci}{\operatorname{Ci}}
\newcommand{\Li}{\operatorname{Li}}
\newcommand{\arctanh}{\operatorname{arctanh}}
\newcommand{\dd}{\,\mathrm{d}}
\newcommand{\QQ}{\mathbb{Q}}
\newcommand{\Qbar}{\overline{\mathbb{Q}}}
\newcommand{\ZZ}{\mathbb{Z}}
\newcommand{\CC}{\mathbb{C}}
\newcommand{\OO}{O}
\newcommand{\Mul}{\mathsf{M}}
\newcommand{\todo}[1]{\textcolor{red}{[TODO: #1]}}

\title[Sums of holonomic functions in subquadratic time]{Computing infinite sums of holonomic functions in subquadratic time}
\author{Fredrik Johansson}
\date{Draft, November 2026}

\raggedbottom
\begin{document}

\begin{abstract}
We show that infinite series $\sum_{k=1}^{\infty} F(k)$, where $F$ is a holonomic function that is regular at
infinity and defined over the algebraic numbers (its differential equation and its expansion at infinity have
algebraic coefficients), can be approximated to $p$ digits in time
$\widetilde O(p^{3/2})$, more precisely with $O(p^{1/2}\Mul(p)\log^2 p)$ bit operations where $\Mul(p)$ is the
cost of $p$-bit integer multiplication. This is a speedup over the classical techniques, such as
Euler--Maclaurin summation and expansion in zeta values, which all cost $\widetilde O(p^2)$ or more. The method
is a generalization of the classical theta-function splitting for the Riemann zeta function: it turns the
sum into a pair of rapidly convergent series of values of holonomic functions with $O(\sqrt p)$ terms each, and
every term is computed in quasilinear time by binary splitting or the bit-burst algorithm. As a result, we obtain the first
subquadratic algorithms for several classical real numbers, including Khinchin's constant, for which we
present a record computation of $10^7$ digits.
\end{abstract}

\maketitle



%======================================================================
\section{Introduction}\label{sec:intro}

A function is \emph{holonomic} (or D-finite) if it satisfies a linear differential equation with polynomial
coefficients, and a sequence is holonomic (or P-recursive) if it satisfies a linear recurrence with polynomial
coefficients \cite{Sta80}. Holonomic functions include the elementary functions, polylogarithms, the
hypergeometric functions and most other special functions of analysis, and their power series coefficients are
holonomic sequences.

Holonomic sequences are behind the fastest known algorithms for many constants. If the terms of a series
$\sum_{k=0}^{\infty} t_k$ form a holonomic sequence with rational coefficients and decrease geometrically, the sum can be
computed to $p$ bits in quasilinear time $O(\Mul(p)\log^2p)$ by binary splitting
(Section~\ref{sec:prelim}), where $\Mul(p)$ denotes the cost of multiplying two $p$-bit integers, which is
$O(p\log p)$ \cite{HvdH21} and even $O(p(\log p)^{1-\kappa})$ for an explicit $\kappa > 0$ \cite{OAI26}. In this way, $e$, $\pi$, $\log 2$, $\zeta(3)$, Catalan's constant, $\Gamma(\frac14)$,
Euler's constant $\gamma$ and many other classical constants can be computed in quasilinear time
\cite{BB87,HP98,BM80,BZ10}, and with the bit-burst algorithm the same holds for the value of any holonomic
function at any computable point \cite{CC90,vdH99}.

\subsection{Sums of holonomic functions}
This paper is about a different class of series: sums
\[
  S = \sum_{k=1}^{\infty} F(k)
\]
in which the summand $F$ is a holonomic function rather than a holonomic sequence. The most prominent
example is Khinchin's constant
\[
  K = \prod_{k=1}^{\infty} \Bigl(1 + \frac{1}{k(k+2)}\Bigr)^{\log_2 k} = 2.6854520010653064453\ldots,
\]
the almost-sure limit of the geometric mean of the partial quotients of a continued fraction
expansion \cite{Khi35}, for which (by summation by parts)
\begin{equation}\label{eq:logK}
  \log K \log 2 = \sum_{k=2}^{\infty} F(k), \qquad F(k) = -\log\Bigl(1-\frac1k\Bigr)\log\Bigl(1+\frac1k\Bigr).
\end{equation}
Such series differ from the classical holonomic series in two ways. Their terms decrease only algebraically,
like $k^{-2}$ here, so that $2^{\Theta(p)}$ terms would be needed for $p$ bits. And the terms $F(k)$ in general do
not form a holonomic sequence (already $\log k$ does not \cite{FGS05}), so neither binary splitting nor the usual
acceleration techniques for holonomic sequences apply to the sum.

\subsection{The class of summands}
We assume that $F$ satisfies a linear differential equation with coefficients in $\Qbar[x]$ for which
$x = \infty$ is an ordinary point or a regular singular point, and that $F$ is \emph{defined over $\Qbar$}: its
coordinates in a Frobenius basis of solutions at infinity are algebraic numbers. Then, for large $x$,
\begin{equation}\label{eq:frob}
  F(x) = \sum_{\sigma \in \Sigma}\sum_{j=0}^{J} x^{-\sigma}(\log x)^j f_{\sigma,j}(1/x)
\end{equation}
with a finite set $\Sigma$ of algebraic exponents and power series $f_{\sigma,j}$ with coefficients in a number
field (Section~\ref{sec:regular}). We also assume that $F(x) = O(x^{-1-\varepsilon})$ as $x \to \infty$, so that
the sum converges absolutely. Some algebraicity condition is needed for a complexity bound, since otherwise
$F$ could be a hard-to-compute constant times a simple function. The regularity condition concerns only the point
at infinity: $F$ may have irregular singularities elsewhere, such as the essential singularity of
$e^{x/k^2}$ at $k = 0$ (an example from plasma physics follows below), but it may not behave exponentially as
$k \to \infty$, which excludes summands such as $z^kF(k)$ with $|z| < 1$ (Section~\ref{sec:zetatheta}).
Regularity at infinity holds for essentially all functions of interest (Remark~\ref{rem:H}).

The core case, which contains all the analysis, is that of $F$ analytic at infinity with a periodic weight:
\begin{equation}\label{eq:S}
  S = \sum_{k=k_0}^{\infty} \chi(k)\, f(1/k), \qquad f(x) = \sum_{n=2}^{\infty} a_n x^n,
\end{equation}
where $f$ is a holonomic power series with algebraic (for simplicity: rational) coefficients and radius of
convergence $\rho > 0$, $\chi$ is periodic, and $k_0 > 1/\rho$ (finitely many terms with $k < k_0$ can be added
separately). For Khinchin's constant, $f(x) = -\log(1-x)\log(1+x)$, $\chi = 1$ and $k_0 = 2$. The general
summands \eqref{eq:frob} reduce to this case by differentiation with respect to the exponent
(Section~\ref{sec:regular}). Periodic weights and factors $(\log k)^j$ do not enlarge the class: a periodic
weight of period $q$ is removed by summing over the residue classes,
\[
  \sum_{k=1}^{\infty} \chi(k)F(k) = \sum_{a=1}^{q}\chi(a)\sum_{j=0}^{\infty} F(qj+a),
\]
and $F(qx+a)$ and $(\log x)^jF(x)$ are again
holonomic, defined over $\Qbar$ and regular at infinity. We keep them in the statements because the splitting
handles them directly.

\subsection{Main result}
We show (Theorem~\ref{thm:general}) that for every $F$ as above and every $k_0$ beyond the finite singular
points of its equation, $\sum_{k=k_0}^{\infty}F(k)$ can be computed to $p$ bits in
\[
  O\bigl(p^{1/2}\, \Mul(p) \log^2 p\bigr) = O\bigl(p^{3/2}(\log p)^{3-\kappa}\bigr)
\]
bit operations (on a multitape Turing machine; Section~\ref{sec:prelim}). The main step is Theorem~\ref{thm:main}, for the sums \eqref{eq:S} with $f$ holonomic and its
differential equation regular singular (or ordinary) at $x = 0$; Theorem~\ref{thm:log} adds powers of
$\log k$, and Theorem~\ref{thm:general} adds non-integer exponents. Previously published methods for general $f$
need $p^{2+o(1)}$ bit operations or more (Section~\ref{sec:related}).

\subsection{The method}\label{sec:method}
The method generalizes Riemann's theta-function splitting for the zeta function, which we recall together with its
other descendants in Section~\ref{sec:zetatheta}. Take $\chi = 1$ for simplicity. Writing each power in
$f(1/k) = \sum_{n=2}^{\infty} a_nk^{-n}$ as a Laplace transform in $k^2$,
\[
  k^{-n} = \frac{1}{\Gamma(n/2)}\int_0^\infty u^{n/2-1}e^{-k^2u}\dd u,
\]
turns the sum into an integral
\[
  S = \int_0^\infty h(u)\,\Theta(u)\dd u, \qquad \Theta(u) = \sum_{k=k_0}^{\infty}e^{-k^2u},
\]
of a theta series against the \emph{kernel}
\[
  h(u) = \sum_{n=2}^{\infty} \frac{a_n u^{n/2-1}}{\Gamma(n/2)},
\]
which is obtained from $f$ by a Borel-type transform and is again holonomic. We split the integral at a point
$u_0 > 0$. On $[u_0, \infty)$ the theta series converges rapidly, and we integrate it term by term: this gives the
\emph{$k$-terms}
\[
  \int_{u_0}^\infty h(u)e^{-k^2u}\dd u,
\]
of which $O(\sqrt p)$ are needed for $p$ bits; for rational $u_0$ they are rational series times $e^{-k^2u_0}$ and
$\erfc(k\sqrt{u_0})$, computed by binary splitting. On $(0, u_0]$ the theta transformation (Poisson summation)
rewrites $\Theta(u)$ as an elementary part ($\frac12\sqrt{\pi/u}$ minus finitely many exponentials), which gives a
\emph{main term}, plus $\sqrt{\pi/u}$ times a series in $e^{-\pi^2m^2/u}$, which gives the \emph{dual terms}, the
moments
\[
  \int_0^{u_0}u^{-1/2}e^{-c/u}h(u)\dd u, \qquad c = \pi^2m^2,
\]
again $O(\sqrt p)$ of them. Integration by parts shows that these moments, as functions of $c$,
satisfy a linear differential system of fixed dimension, which we call the \emph{dual system}; each dual term is
computed from it with the bit-burst algorithm. Every term costs $O(\Mul(p)\log^2 p)$, which gives the bound above.

The algebraicity assumption matters only for the complexity. If the coefficients of the differential equation or of the expansion \eqref{eq:frob} are merely computable real or
complex numbers, for instance because they depend on a parameter given only by its approximations, the method still
applies as a convergence acceleration technique: the splitting and all identities are unchanged, and $O(\sqrt p)$
terms on each side still suffice. But the series for the individual terms then have coefficients that are
full-precision numbers rather than small rationals, so that binary splitting no longer applies; with baby-step
giant-step evaluation of holonomic sequences \cite{BGS07,Joh14}, each term costs $p^{3/2+o(1)}$ instead of
$p^{1+o(1)}$, and the total complexity degenerates to $p^{2+o(1)}$, like that of the classical methods.

\subsection{New examples}
The class \eqref{eq:S} contains many constants for which, to our knowledge, only quadratic or slower
algorithms were known before:
\begin{itemize}
\item Khinchin's constant $K$ (Section~\ref{sec:khinchin}) and its H\"older means $K_{-n}$, $n \ge 1$, such as the
harmonic mean $K_{-1}$ (Section~\ref{sec:KH});
\item $\mathcal C = \sum_{k=1}^{\infty}\frac{\log k}{k(k+1)}$, which gives Kalpazidou's analogue $e^{\mathcal C}$ of Khinchin's constant
for L\"uroth expansions and the Alladi--Grinstead constant $e^{\mathcal C - 1}$ (Section~\ref{sec:AG});
\item $\mathcal A = \sum_{k=2}^{\infty}\frac{\log k}{k^2-1}$, related to L\"uroth expansions (Section~\ref{sec:AG});
\item sums of inverse trigonometric functions recorded in the OEIS \cite{OEIS} to about $90$ digits, such as
(A323736, A323737, A343470)
\begin{gather*}
  \lim_{k\to\infty}\Bigl(\sum_{j=1}^{k}\arcsin\frac1j - \log k\Bigr)
    = \gamma + \sum_{j=1}^{\infty}\Bigl(\arcsin\frac1j - \frac1j\Bigr), \\
  \sum_{j=1}^{\infty}\Bigl(j\arcsin\frac1j - 1\Bigr), \\
  \sum_{n=1}^{\infty}(-1)^{n-1}\frac{\arctan n}{n} = \frac\pi2\log2 - \sum_{n=1}^{\infty}\frac{(-1)^{n-1}}{n}\arctan\frac1n;
\end{gather*}
\item $\sum_{k=2}^{\infty}\zeta''(k) = \sum_{k=2}^{\infty}\frac{\log^2 k}{k(k-1)}$ \cite[p.~538]{Fin18} (A345385), and more generally sums
$\sum_{k=1}^{\infty} (\log k)^j R(k)$ with rational $R$, such as $\sum_{k=1}^{\infty} \frac{\log k}{k^2+1}$;
\item sums of products of logarithms, polylogarithms and inverse trigonometric functions of $1/k$, such as
$\sum_{k=1}^{\infty} \frac1k\arctan\frac1k$ and Boyadzhiev's constant $\sum_{k=1}^{\infty}\frac1k\log(1+\frac1k)$ \cite{Boy19}, and products
$\prod_{k=1}^{\infty} G(1/k)$ with $\log G$ holonomic;
\item sums with elliptic-integral summands from physics, such as the velocity induced in a periodic array of
coaxial vortex rings with spacing $\lambda$ \cite{MS22}, whose $n$-th term is a combination of complete elliptic
integrals $K$ and $E$ of modulus $4/(4+n^2\lambda^2)$ (Section~\ref{sec:vortex});
\item summands with an essential singularity at $k = 0$ (but analytic at infinity), such as the
Brillouin--Planck--Larkin partition function of the hydrogen atom,
\[
  \sigma(T) = \sum_{s=1}^{\infty}s^2\Bigl(e^{x/s^2} - 1 - \frac{x}{s^2}\Bigr), \qquad x = \frac{I}{k_BT},
\]
the subtracted form of the divergent
sum over the Bohr levels used in plasma physics \cite{EKR11}, and, at rational $x$, the Hardy--Littlewood
function $\sum_{k=1}^{\infty}\frac1k\sin\frac xk$ \cite{Gau05}. Levels with quantum defects $\delta$ give summands
with an essential singularity at $s = \delta$.
\end{itemize}

\subsection{Computations}
\subsubsection{Khinchin's constant}
Applied directly to $f(x) = -\log(1-x)\log(1+x)$, the method gives a dual system of order six.
A dilogarithm identity of Bailey, Borwein and Crandall (\cite[\S3]{BBC97}; Lemma~\ref{lem:dilog}) rewrites
$\log K\log 2$ as an alternating sum of $\Li_2(4/k^2)$, for which the kernel is an exponential integral, $\Eip(4u)/u$ with $\Eip(y) = \sum_{n=1}^{\infty}y^n/(n\,n!)$,
the dual terms have a closed form in terms of $\erfc$, the cosine integral and an entire power series
(Theorem~\ref{thm:khinchin}), and the dual system has order four with coefficients linear in $c$. With this
formula, rigorous error bounds, and an implementation in FLINT \cite{FLINT}, we have computed $K$ to $10^7$ digits in $11.4$ hours on $8$ threads, against
$65$ hours for the zeta-value expansion in FLINT, which we ran for verification (Section~\ref{sec:khresults}). The
digits determine the first $9\,699\,142$ partial quotients of the continued fraction of $K$, and show that $K$ is not a
rational number with denominator below $2.58\cdot10^{5\,000\,000}$ (Theorem~\ref{thm:khrat}). The digits of $K$
have been computed to $7\,350$ digits by Bailey, Borwein and Crandall \cite{BBC97}, to
$110\,000$ digits by Gourdon (1997), to $10^6$ digits by Sim\'o (2016) and to $2\cdot10^6$ digits by
Ghazi (2026) \cite{Gha26}, according to \cite{MWKdigits}. All of these used the zeta-value expansion, and all previous
algorithms take at least quadratic time; in the words of Weisstein \cite{MWKdigits}, ``the numerical value of $K$ is
notoriously difficult to calculate to high precision''.

\subsubsection{Other constants}
Section~\ref{sec:apps} applies the method to a few further constants that,
to our knowledge, had only been computed to a few hundred or thousand digits: the constants $\mathcal A$ and
$\mathcal C$ above and the harmonic mean $K_{-1}$, as well as the vortex sum and the Hardy--Littlewood function. We
compute each of them to $10^6$ digits, and compare with the zeta-value expansion and Euler--Maclaurin summation up to
$10^5$ digits.

%----------------------------------------------------------------------
\subsection{Zeta series and theta splitting}\label{sec:zetatheta}

\subsubsection{Riemann's splitting}
The simplest summand is a pure power $F(k) = k^{-s}$, and $S = \zeta(s)$; with periodic weights and shifts one
obtains Dirichlet $L$-functions and the Hurwitz and Lerch zeta functions. We call these \emph{zeta series}.
Riemann's second proof of the functional equation \cite{Rie59} writes
\[
  \pi^{-s/2}\Gamma(s/2)\zeta(s) = \int_0^\infty u^{s/2-1}\,\frac{\theta(u)-1}{2}\dd u, \qquad
  \theta(u) = \sum_{k\in\ZZ} e^{-\pi k^2 u},
\]
splits the integral at $u = 1$ and transforms the part over $(0, 1]$
with $\theta(u) = u^{-1/2}\theta(1/u)$. The result is a series for $\zeta(s)$ in incomplete gamma functions that
converges geometrically on both sides. Its generalizations are the approximate functional equations of Dirichlet and
more general $L$-functions \cite{Lav68,CN62}, the basis of the numerical methods of Lagarias and Odlyzko,
Rubinstein \cite{Rub05}, Dokchitser \cite{Dok04} and others, and of the approximate functional equations of the
Lerch zeta function \cite{GLS03}. In mathematical physics the same device is Ewald summation of lattice sums
\cite{Ewa21}, the Chowla--Selberg formula for Epstein zeta functions \cite{CS67} (see \cite{BGMWZ13} for an extensive
account), and the proper-time and heat-kernel representations of $\sum_{k=1}^{\infty} F(k^2)$ in quantum field theory
\cite{Sch51,Eli95,Kir02}. Crandall developed these splittings systematically as computational tools
(``Riemann splitting''): for $\zeta$ and related functions in \cite{BBC00}, for Madelung-type and Epstein sums, and
for polylogarithms, Lerch transcendents and $L$-series in \cite{Cra12} (see also \cite{Cof12}).

\subsubsection{Complexity for zeta series}
At integers $s$, zeta values can be computed in quasilinear time \cite{Kar95,Kar98}, by binary splitting of rapidly
convergent series. For algebraic $s$ of bounded degree, the approximate functional equation of $L(s,\chi)$ combined
with fast evaluation of the incomplete gamma function gives $p^{3/2+o(1)}$ \cite{Joh21}; this extends to
derivatives with respect to $s$, that is, to sums $\sum_{k=1}^{\infty} \chi(k)(\log k)^j k^{-s}$, and gives for example the
Stieltjes constants (see also \cite{Cof14}) and Glaisher's constant $A$, through $\log A = \frac1{12} - \zeta'(-1)$
(the method used in FLINT), in time $p^{3/2+o(1)}$. For a generic complex $s$, given only by approximations, the series for the terms can no longer be
summed by binary splitting, and the cost becomes $p^{2+o(1)}$.

\subsubsection{Example: the constants of Euler and Glaisher}
The Laurent expansion of Riemann's formula at $s = 1$ gives, with $\Gamma(\frac12, x) = \sqrt\pi\erfc\sqrt x$ and
$\Gamma(0, x) = E_1(x)$ (Coffey \cite[Proposition~9]{Cof14}),
\begin{equation}\label{eq:gammaRiemann}
  \gamma = 2\sum_{k=1}^{\infty}\Bigl(\frac{\erfc(k\sqrt\pi)}{k} + E_1(\pi k^2)\Bigr) - 2 + \log(4\pi),
\end{equation}
and the next coefficient gives a similar but longer series for the Stieltjes constant $\gamma_1$ \cite[(1.22)]{Cof14}.
These series are implicit in \cite{Rie59}. In the terms of this paper, they result from applying the splitting
to the divergent sum $\sum_{k=1}^{\infty} 1/k$ ($f(x) = x$, at the self-dual point $u_0 = \pi$): only the main term diverges, and
taking its finite part by analytic continuation in the exponent ($k^{-s}$ with $s \to 1$) gives \eqref{eq:gammaRiemann},
while the derivatives in $s$, as in Theorem~\ref{thm:log}, give $\gamma_1$ and the higher Stieltjes constants. Any
other $u_0$ gives a formula of the same shape. The kernel is the monomial $(\pi u)^{-1/2}$, so that every term is a
single incomplete gamma function, with $k$-terms and dual terms of the same size $e^{-\pi k^2}$. Euler's constant is also the sum of a holonomic function,
\[
  \gamma = \sum_{k=1}^{\infty}\Bigl(\frac1k - \log\Bigl(1+\frac1k\Bigr)\Bigr), \qquad f(x) = x - \log(1+x) = \sum_{n=2}^{\infty}\frac{(-1)^nx^n}{n},
\]
to which Theorem~\ref{thm:main} applies, but it produces a different and more complicated series: $f$ has radius of
convergence $1$, so the term $k = 1$ is taken out; the kernel is
\[
  h(u) = \frac{e^u-1}{2u} - \sum_{m=1}^{\infty}\frac{u^{m-1/2}}{(2m+1)\Gamma(m+\frac12)};
\]
and every $k$-term and dual term is a series rather than a single incomplete gamma function. (The method is linear
in $f$, and the regularized contributions of the $x^1$ terms of $x$ and $\log(1+x)$ cancel.) Likewise
\[
  \gamma_1 = \sum_{k=1}^{\infty}\Bigl(\log k\,\Bigl(\frac1k - \log\Bigl(1+\frac1k\Bigr)\Bigr)
    - \frac12\log^2\Bigl(1+\frac1k\Bigr)\Bigr)
\]
is covered by Theorems~\ref{thm:main} and~\ref{thm:log}, but the result is not Coffey's series. Of course, $\gamma$ can be computed in
quasilinear time \cite{BM80}, and $\gamma_1$ and $\zeta'(-1)$ in time $p^{3/2+o(1)}$ by \cite{Joh21}; we make no
claim to improve on either, and give these examples only to illustrate the shape of the method. In the same way,
Glaisher's constant is a sum of a holonomic
function without logarithmic weights:
\begin{align*}
  \sum_{k=1}^{\infty}\Bigl(3k + \frac12 - k(3k+2)\log\Bigl(1+\frac1k\Bigr)\Bigr)
  &= \sum_{n=2}^{\infty}\frac{(-1)^n(n-1)}{(n+1)(n+2)}\zeta(n) \\
  &= 6\log A - \frac12 - \frac12\log(2\pi),
\end{align*}
with $f(x) = 3/x + \frac12 - (3+2x)x^{-2}\log(1+x)$ (the first equality is the zeta-value expansion, the second
follows from the asymptotics of the hyperfactorial $\prod_{k=1}^{N}k^k$).

\subsubsection{Beyond pure powers}
In almost all of the works above the summand is a power $k^{-s}$, or a power of a quadratic form, so that the kernel
is $u^{s-1}$ and every term of the splitting is a single incomplete gamma function. A notable exception are the
inhomogeneous sums $\sum_{k=1}^{\infty} (k^2+a^2)^{-s}$ of the Chowla--Selberg formula and of zeta regularization
\cite{CS67,Eli95,Kir02}, whose kernel is $u^{s-1}e^{-a^2u}$ and whose dual terms are modified Bessel functions; in our
language this is $f(x) = x^{2s}(1+a^2x^2)^{-s}$, a holonomic kernel with dual terms that are values of a holonomic
function. We show that this works for every summand in our class: the kernel is always holonomic, and the dual terms
are always values of a holonomic function of $c$ with no singularity in $\operatorname{Re} c > 0$
(Lemma~\ref{lem:moments}). Theorem~\ref{thm:main} contains the $L$-function result of \cite{Joh21} at integer $s$ as
the special case $f(x) = x^s$ (where the kernel is a monomial), and Theorem~\ref{thm:general} contains it at algebraic
$s$, but not its uniformity in the height of $s$, which is what gives the Bernoulli number results; see
Remark~\ref{rem:uniform}.

\subsubsection{Dirichlet series}
The same splitting applies to the Dirichlet series $D_F(s) = \sum_{k=1}^{\infty}F(k)k^{-s}$, since $F(k)k^{-s}$ is again
in our class for algebraic $s$. For $F(k) = f(1/k)$ as in \eqref{eq:S}, the expansion
\[
  D_F(s) = \sum_{k=1}^{k_0-1}F(k)k^{-s} + \sum_{n=2}^{\infty}a_n\zeta(s+n, k_0)
\]
converges for all $s$ away from the poles, so $D_F$ is meromorphic in $\CC$ with at most simple poles at $s = 1-n$, of residue $a_n$, and polynomial growth in vertical
strips; factors $k^{-\sigma}(\log k)^j$ move the poles to $1 - \sigma - n$ and raise their order to $j+1$. In the
splitting, the $k$-side and dual terms are entire in $s$ and the poles come from the main term, so the splitting
formula is itself the continuation, as in Riemann's proof. In particular the zeta-regularized values of divergent
sums $\sum_{k=1}^{\infty} k^mF(k)$, which occur in Casimir-type calculations, are computed in the same time. Dirichlet series
attached to polynomials and rational functions, $\sum_{k=1}^{\infty} P(k)^{-s}$, are in the class for every fixed algebraic $s$;
their meromorphic continuation and special values have been studied since Mellin and Mahler \cite{Mah28}, for instance
by Cassou-Nogu\`es \cite{CN81}, Sargos \cite{Sar84} and Peter \cite{Pet98}, but not their computation. What is lost in
general is the functional equation: Poisson summation relates $D_F(s)$ to a series of values of a holonomic function
at the nonzero integers (our dual terms), which is not a Dirichlet series except for special $F$ such as powers
and the Chowla--Selberg summands. Approximate functional equations of general Dirichlet series are derived from
a functional equation \cite{Lav68,CN62}, and without some structure no acceleration is possible: for a bounded random
sequence $a_k$, the sum $\sum_{k=1}^{\infty}a_kk^{-2}$ needs $2^{\Theta(p)}$ terms for $p$ bits. The two known sources of
acceleration are a functional equation (automorphy) and smoothness of $k \mapsto a_k$, which makes Poisson summation
effective; our class is a systematic instance of the second.

\subsubsection{Holonomic sequences}
Dirichlet series whose coefficients form a holonomic sequence, such as the harmonic zeta function
$\sum_{k=1}^{\infty} H_kk^{-s}$ \cite{AV84,CC20}, are a different class. Holonomic sequences of polynomial growth typically have
divergent asymptotic expansions: $H_k = \psi(k+1)+\gamma$ and ratios of gamma functions such as $\binom{2k}{k}4^{-k}$
are not values of a function that is holonomic and regular at infinity, and are not covered by our method.
Conversely, $\log k$ and $\sqrt k$ are not holonomic sequences \cite{FGS05}. The overlap of the two classes is
small; it contains the rational functions of $k$.

\subsubsection{Reductions to zeta series}
Some sums with holonomic summands reduce to zeta series. If $Q$ is a rational function whose poles all lie
at the origin, then $\sum_{k=1}^{\infty} \chi(k)(\log k)^jQ(k)$ is a finite combination of derivatives of Dirichlet $L$-functions
or Lerch zeta functions at integers; likewise, replacing $\log k$ by $\log(k+a)$ and the poles by poles at $-a$ gives
derivatives of Hurwitz zeta values. In general, partial fractions produce terms
\[
  \frac{\log k}{(k+a)^m} = \frac{\log(k+a)}{(k+a)^m} - \frac{\log(1+a/k)}{(k+a)^m},
\]
and the second part is again a sum of the form \eqref{eq:S}, so this only moves the problem: an example is
\[
  \mathcal C = \sum_{k=1}^{\infty}\frac{\log k}{k(k+1)} = -\sum_{k=2}^{\infty}\frac1k\log\Bigl(1-\frac1k\Bigr).
\]
Khinchin's constant is of
the form $\sum_{k=2}^{\infty} \log k\, g(1/k)$ with $g$ analytic but not rational,
\[
  \log K\log 2 = \sum_{k=2}^{\infty}\log k\,\Bigl(2\log\Bigl(1+\frac1k\Bigr) - \log\Bigl(1+\frac2k\Bigr)\Bigr),
\]
and expanding $g$ only gives an infinite series in the values $\zeta'(n)$. No reduction of $K$, $\mathcal C$
or the other examples below to finitely many zeta series values is known.

\subsubsection{Exponential weights}
Weights $e^{2\pi i\alpha k}$ with real algebraic $\alpha$, as in the Lerch zeta function, can be handled by a theta
function with characteristic (Remark~\ref{rem:ext}). Weights $z^k$ with $|z| < 1$ cannot: then $z^kF(k)$ has an
irregular singularity at infinity, the even extension used in the splitting does not exist, and the one-sided
sums $\sum_{k=1}^{\infty}z^ke^{-k^2u}$ are partial theta functions, which have no modular transformation. Crandall's
algorithms for the Lerch transcendent and polylogarithms \cite{Cra12} cover all $z$; for $|z| < 1$, splitting the
Laplace representation in $k$ (rather than $k^2$) gives two geometrically convergent series, but with $\Theta(p)$
terms. These sums converge geometrically, but
their terms are not holonomic sequences, and we do not know a subquadratic algorithm even for the polylogarithm
$\Li_s(z)$ with rational $z \in (0, 1)$ and algebraic non-integer $s$.

%----------------------------------------------------------------------
\subsection{Other methods}\label{sec:related}

\subsubsection{The zeta-value expansion}
Expanding $f$ in \eqref{eq:S} and interchanging the sums gives
\begin{equation}\label{eq:zetaseries}
  S = \sum_{n=2}^{\infty} a_n \zeta_{\chi, k_0}(n), \qquad \zeta_{\chi,k_0}(n) = \sum_{k=k_0}^{\infty} \chi(k) k^{-n},
\end{equation}
a combination of values of Hurwitz zeta functions or Dirichlet $L$-functions at the integers. For $\log K\log 2$ this
is the series
\[
  \log K\log 2 = \sum_{n=1}^{\infty} \frac{\zeta(2n)-1}{n} \Bigl(1 - \frac12 + \cdots + \frac{1}{2n-1}\Bigr)
\]
of Shanks and Wrench \cite{SW59}, used with refinements by Bailey, Borwein and Crandall \cite{BBC97}, and Flajolet and
Vardi \cite{FV96} list many further constants with such expansions. We call \eqref{eq:zetaseries} the
\emph{zeta-value expansion}. Its terms decay like $(k_0\rho)^{-n}$, so $p$ bits require about $p/\log_2(k_0\rho)$
zeta values to $p$ bits each. Their total size alone is of order $p^2$, so every algorithm that goes through
\eqref{eq:zetaseries} needs $p^{2+o(1)}$ bit operations, and the practical algorithms (power sums, Bernoulli
numbers, Euler products for the individual $\zeta(n)$) need more.

\subsubsection{Summation and integration methods}
Direct summation of \eqref{eq:S} with the Euler--Maclaurin formula or the Abel--Plana formula, Gaussian quadrature
for sums \cite{Mon10} or with Bose--Einstein weights \cite{GM85}, and numerical integration of a Mellin--Barnes or
contour integral representation all need $\Theta(p)$ function evaluations or terms at precision $p$, hence again at
least $p^{2+o(1)}$; see Belabas and Cohen \cite{BC21} for a survey of these methods and their implementations in
PARI/GP, and Flajolet, Gourdon and Dumas \cite{FGD95} for the related Mellin transform asymptotics of harmonic sums
$\sum_{k=1}^{\infty} F(hk)$. Convergence acceleration of the Cohen--Rodriguez Villegas--Zagier type \cite{CRVZ00} does not apply
directly (the sums are not alternating), and where it does, it needs $\Theta(p)$ terms as well. To our knowledge no
subquadratic algorithm for \eqref{eq:S} with general holonomic $f$ has been published. Apart from the zeta series, some
special instances are classical constants with quasilinear algorithms, such as Euler's constant (Brent and
McMillan \cite{BM80}), $\sum_{k=2}^{\infty}(\frac1k - \log(1+\frac1k)) = \gamma - 1 + \log 2$.

\subsubsection{Fast evaluation of the pieces}
Binary splitting evaluates convergent series with rational P-recursive terms to $p$ bits in
$O(\Mul(p)\log^2 p)$ \cite{HP98,CHTZZ07}; the bit-burst algorithm of Chudnovsky and
Chudnovsky \cite{CC90} and van der Hoeven \cite{vdH99} extends this to values of holonomic
functions at arbitrary computable points; see also Mezzarobba \cite{Mez10,Mez16} for rigorous
implementations.

%----------------------------------------------------------------------
\subsection{Outline}
Section~\ref{sec:prelim} fixes the model of computation and recalls binary splitting and the
bit-burst algorithm. Section~\ref{sec:general} proves the main theorems, Section~\ref{sec:khinchin} treats
Khinchin's constant in detail, including our record computation and its consequences for the continued fraction of
$K$, and Section~\ref{sec:apps} the other constants.

%======================================================================
\section{Preliminaries}\label{sec:prelim}

\subsection{Model of computation}
Throughout, ``time'' means bit complexity on a multitape Turing machine. We let $\Mul(p)$ be a function such that
two $p$-bit integers can be multiplied with $O(\Mul(p))$ bit operations in this model, and assume the usual
regularity conditions
\begin{equation}\label{eq:Mreg}
  \Mul(p)/p \text{ is nondecreasing}, \qquad \Mul(p\log p) = O(\Mul(p)\log p).
\end{equation}
One can take $\Mul(p) = p\log p$ \cite{HvdH21}, and even $\Mul(p) = p(\log p)^{1-\kappa}$ with $\kappa = 2^{-182}$
\cite{OAI26} (larger values of $\kappa$ have since been obtained); both satisfy \eqref{eq:Mreg}. The first condition
gives $\Mul(a) + \Mul(b) \le \Mul(a+b)$, so that each level of a product tree costs at most one multiplication of
the total size, and geometric sums of multiplication costs, as in Newton iteration, are $O(\Mul(p))$. Together with
the first, the second gives $\Mul(Cp) = O(\Mul(p))$ for constant $C$, and it absorbs the $O(\log p)$-bit
coefficients in binary splitting (Sections~\ref{sec:bsplit} and~\ref{sec:bitburst}). No other property of $\Mul$
is used, so all bounds below hold for any such $\Mul$, with no additional logarithmic factors when
$\Mul(p) = o(p\log p)$. Besides integer
multiplication, the algorithm uses only binary splitting, the bit-burst algorithm, and division, square
roots and exponentials by Newton iteration and binary splitting, whose bit complexity bounds hold in this
model \cite{BZ10} (for Theorem~\ref{thm:general} also logarithms, and the gamma function at algebraic
points, by the same methods); everything else is a sequence of $O(\sqrt p)$ arithmetic operations on numbers of $O(p)$
bits and bookkeeping with integers of $O(\log p)$ bits.

\subsection{Notation}
We write $\widetilde O(g) = g\,(\log g)^{O(1)}$ and $\vartheta = u\,\frac{\mathrm{d}}{\mathrm{d}u}$ (or
$x\frac{\mathrm{d}}{\mathrm{d}x}$), and $(a)_n = a(a+1)\cdots(a+n-1)$. We use the truncated exponential
$e_n(x) = \sum_{i=0}^{n-1} x^i/i!$ (note: $n$ terms) and the incomplete gamma and exponential integral functions
\begin{gather*}
  \Gamma(s, y) = \int_y^\infty t^{s-1}e^{-t}\dd t, \qquad Q(s,y) = \frac{\Gamma(s,y)}{\Gamma(s)}, \\
  E_1(y) = \Gamma(0, y) = -\gamma - \log y + \Ein(y), \qquad \Ein(y) = \sum_{n=1}^{\infty}\frac{(-1)^{n-1}y^n}{n\,n!}, \\
  \Eip(y) = -\Ein(-y) = \sum_{n=1}^{\infty}\frac{y^n}{n\,n!} = \operatorname{Ei}(y) - \gamma - \log y,
\end{gather*}
so that $Q(m, y) = e^{-y} e_m(y)$ and $\Gamma(\frac12, y) = \sqrt\pi\erfc\sqrt y$. ``Computing $S$ to $p$ bits''
means producing an approximation with absolute error at most $2^{-p}$.

\subsection{Binary splitting}\label{sec:bsplit}
Let
\[
  X_{n+1} = q_n^{-1} A_n X_n, \qquad s_{n} = s_{n-1} + q_n^{-1} B_n X_n,
\]
where $A_n$, $B_n$ are matrices of fixed size whose entries, like $q_n$, are integers of $O(\log p)$ bits for
$n \le N = O(p)$. Then $X_N$ and $s_N$ are rational numbers whose numerators and denominators can be computed by
the product tree of the matrices $\bigl[\begin{smallmatrix} A_n & 0 \\ B_n & q_n\end{smallmatrix}\bigr]$
\cite{HP98}. Each of its $O(\log p)$ levels has entries of total size $O(p\log p)$, so by \eqref{eq:Mreg} the cost is
\[
  O(\Mul(p\log p)\log p) = O(\Mul(p)\log^2 p).
\] We
use this for series whose terms are P-recursive with rational coefficients and decay at least
geometrically: then $O(p)$ terms suffice for $p$ bits and the result costs $O(\Mul(p) \log^2 p)$.
This covers, at rational points inside the disk of convergence, all holonomic functions with rational
power series coefficients, and the functions $\exp$, $\log$, $\erf$, $\Gamma(\frac12, y)$, $E_1(y)$
at rational arguments $y = O(p)$, using the lower series
\[
  \gamma(\tfrac12,y) = y^{1/2}e^{-y}\sum_{j=0}^{\infty} \frac{y^j}{(\frac12)_{j+1}}
\]
and $E_1(y) = -\gamma - \log y + \Ein(y)$, with $O(p)$ guard bits for the cancellation. Euler's constant $\gamma$
costs $O(\Mul(p)\log^2p)$ \cite{BM80}, and $\pi$ costs $O(\Mul(p)\log p)$.

\subsection{Bit-burst}\label{sec:bitburst} Let $Y(c) \in \CC^D$ satisfy a linear system
\begin{equation}\label{eq:fuchsys}
  c\, Y'(c) = (B_0 + c B_1)\, Y(c), \qquad B_0, B_1 \in \QQ^{D \times D},
\end{equation}
whose only finite singular point is $c = 0$. Its Taylor coefficients at $c_0 \ne 0$,
$Y(c_0 + \delta) = \sum_{l=0}^{\infty} Y_l \delta^l$, satisfy
\begin{equation}\label{eq:taylorrec}
  c_0 (l+1) Y_{l+1} = (B_0 + c_0 B_1 - l) Y_l + B_1 Y_{l-1},
\end{equation}
a recurrence with rational coefficients if $c_0 \in \QQ$, and the series converges for $|\delta| < |c_0|$.
Given $Y(c_0)$ at a rational point $c_0$ and a real number $c$ with $|c - c_0| \le |c_0|/2$, whose
approximations to $p$ bits can be computed in $O(\Mul(p)\log^2 p)$, the bit-burst algorithm
\cite{CC90,vdH99} computes $Y(c)$ to $p$ bits in $O(\Mul(p)\log^2 p)$: one chooses rational
$c_j$ with $2^j$-bit binary expansions converging to $c$, and steps from $c_{j}$ to $c_{j+1}$ by
summing \eqref{eq:taylorrec} with binary splitting. The $j$-th step has $O(p/2^j)$ terms with
$O(2^j + \log p)$-bit coefficients, so each level of its product tree has total size $O(p + 2^{-j}p\log p)$,
and the step costs $O(\Mul(p + 2^{-j}p\log p)\log p)$. The steps with $2^j \ge \log p$ cost $O(\Mul(p)\log p)$
each, and there are $O(\log p)$ of them; by \eqref{eq:Mreg} the others form a geometric sum bounded by
$O(\Mul(p\log p)\log p) = O(\Mul(p)\log^2 p)$.
The constants depend on $D$, on the bit size of $B_0, B_1$, on $|c_0|$ and on an a priori bound
$2^{O(p)}$ for $|Y|$ on the disk $|c - c_0| \le \frac34 |c_0|$.

\subsection{Holonomic functions and the condition at the origin}
A power series is holonomic (D-finite) if it satisfies a linear differential equation with polynomial coefficients.
Every such equation can be written in the form
\begin{equation}\label{eq:thetaform}
  \sum_{j=0}^{s} x^j q_j(\vartheta)\, y = 0, \qquad q_j \in \QQ[X],\ q_0 \ne 0 .
\end{equation}
\begin{definition}
The operator \eqref{eq:thetaform} satisfies condition \textbf{(H)} if $\deg q_0 \ge \deg q_j$ for all $j$.
\end{definition}
By Fuchs's criterion, (H) holds if and only if $x = 0$ is an ordinary point or a regular singular
point of the equation, i.e.\ all its local solutions at $0$ have moderate growth \cite[\S15.3]{Inc26}.
In terms of the coefficients, \eqref{eq:thetaform} means $\sum_{j=0}^{s} q_j(n-j)\, a_{n-j} = 0$ for $n$ large,
and (H) says that $q_0$, the coefficient of the newest term, has maximal degree.

\begin{remark}\label{rem:H}
Condition (H) is a condition at $x = 0$, that is, at $k = \infty$ for the summand $f(1/k)$. It holds for all
G-functions (by the theorem of Chudnovsky and Andr\'e their minimal operators are Fuchsian), hence for algebraic
functions, logarithms, polylogarithms and ${}_{q+1}F_q$ functions with rational parameters, and also for entire
functions such as $e^x$, $\erf$ or ${}_pF_q$ with $p \le q$: these have their irregular singularity at $x = \infty$,
which corresponds to $k = 0$ and is never reached by the sum (for instance $f(x) = x^{-2}(e^{cx^2} - 1 - cx^2)$
gives the Planck--Larkin summand). What (H) excludes is exponential behaviour of the summand as $k \to \infty$,
such as a factor $z^k$ with $|z| < 1$. It is preserved
by sums, products, $x \mapsto -x$, $f \mapsto x f$ and $f \mapsto f(x^2)$, since all of these preserve
moderate growth of the solution spaces. A holonomic power series with positive radius of convergence
violating (H) for every annihilating operator would need an irregular singularity at $0$ in all of
them; we do not know whether such series exist.
\end{remark}

%======================================================================
\section{The general theorem}\label{sec:general}

\begin{theorem}\label{thm:main}
Let $f(x) = \sum_{n=2}^{\infty} a_n x^n \in \QQ[[x]]$ be holonomic, annihilated by an operator satisfying (H),
with radius of convergence $\rho > 0$. Let $q \ge 1$, let $\chi \colon \ZZ \to \QQ(\zeta_q)$ have period $q$, and
let $k_0 > 1/\rho$ be an integer. Then
\[
  S = \sum_{k=k_0}^{\infty} \chi(k) f(1/k)
\]
can be computed to $p$ bits in $O(p^{1/2}\Mul(p)\log^2 p)$ bit operations.
The implied constant depends on $f$, $\chi$ and $k_0$.
\end{theorem}

The rest of this section proves the theorem. We fix $u_0 \in \QQ_{>0}$ (independent of $p$) and a
number $R$ with $1/\rho < R < k_0$, so that $|a_n| \le C_R R^n$ for all $n$.

\subsection{Reduction to even weights}
Write $\chi = \chi_0 + \chi_1$ with
\[
  \chi_\delta(k) = \tfrac12\bigl(\chi(k) + (-1)^\delta\chi(-k)\bigr);
\]
both parts are $q$-periodic and $\chi_\delta(-k) = (-1)^\delta \chi_\delta(k)$. Since
$\chi_\delta(k) f(1/k) = w_\delta(k)\, f^{[\delta]}(1/k)$ with
\[
  w_\delta(k) = \chi_\delta(k)\, k^\delta, \qquad f^{[\delta]}(x) = x^\delta f(x),
\]
and $x^\delta f$ again satisfies the hypotheses, it suffices to treat a single $\delta \in \{0,1\}$, a
$q$-periodic $\chi$ of parity $\delta$, and the \emph{even} weight $w(k) = \chi(k)k^\delta$, with
$f = f^{[\delta]}$, whose coefficients we still denote by $a_n$ (now $a_n = 0$ for $n < 2 + \delta$).

\subsection{The kernel}
Write $f(x) = A_0(x^2) + x A_1(x^2)$ and define
\begin{equation}\label{eq:kernels}
\begin{gathered}
  h_0(u) = \sum_{m=1}^{\infty} a_{2m}\frac{u^{m-1}}{\Gamma(m)}, \qquad
  h_1(u) = \sum_{m=1}^{\infty} a_{2m+1} \frac{u^{m}}{\Gamma(m+\frac12)}, \\
  h(u) = h_0(u) + u^{-1/2} h_1(u).
\end{gathered}
\end{equation}

\begin{lemma}\label{lem:kernel}
\begin{enumerate}
\item $h_0$ and $h_1$ are entire, and
\[
  |h_0(u)| \le C R^2 e^{R^2|u|}, \qquad |h_1(u)| \le C (1 + |u|)e^{R^2|u|}.
\]
\item For every integer $k \ge k_0$, $f(1/k) = \int_0^\infty e^{-k^2 u} h(u)\dd u$, with absolute convergence.
\item For $\varepsilon \in \{0,1\}$ there are $r \ge 1$, a matrix $A_\varepsilon(u) \in \QQ[u]^{r\times r}$ and a vector
$b_\varepsilon(u) \in \QQ[u]^r$ such that the vector
\[
  G_\varepsilon = (h_\varepsilon, \vartheta h_\varepsilon, \dots, \vartheta^{r-1}h_\varepsilon)^T
\]
satisfies
\begin{equation}\label{eq:Gsys}
  \vartheta G_\varepsilon = A_\varepsilon(u)\, G_\varepsilon + b_\varepsilon(u).
\end{equation}
\end{enumerate}
\end{lemma}

\begin{proof}
(1) Termwise, from $|a_n| \le C_R R^n$ and
\[
  \sum_{m=0}^{\infty} \frac{x^m}{\Gamma(m+\frac12)} \le (1+\sqrt x)e^x \qquad (x \ge 0).
\]
(2) For $n \ge 2$,
\[
  k^{-n} = \frac{1}{\Gamma(n/2)}\int_0^\infty u^{n/2-1}e^{-k^2u}\dd u;
\]
sum over $n$ and interchange, which is justified by Tonelli's theorem since $\sum_{n=2}^{\infty} |a_n| k^{-n} < \infty$ for
$k > 1/\rho$.

(3) First, $A_0$ and $A_1$ are annihilated by operators satisfying (H). Let $L$ be the operator for $f$ and
$L^-$ the operator obtained by $x \mapsto -x$ (it also satisfies (H)). The least common left multiple
of $L$ and $L^-$ annihilates $f(x) \pm f(-x)$, and its solution space is the sum of the solution spaces of $L$
and $L^-$; by Fuchs's criterion it is regular singular at $0$, i.e.\ satisfies (H). So the even part
$f_e(x) = A_0(x^2)$ is annihilated by an operator $\sum_{j=0}^{s} x^j q_j(\vartheta)$ with (H). Split it as
$L_{\mathrm{ev}} + L_{\mathrm{od}}$ where $L_{\mathrm{ev}}$ collects the even powers $x^j$; since $\vartheta$ preserves parity and $f_e$ is even,
$L_{\mathrm{ev}} f_e$ is even and $L_{\mathrm{od}} f_e$ is odd, so both vanish. With $y = x^2$, $\vartheta_x = 2\vartheta_y$, this reads
$\sum_{i=0}^{\lfloor s/2\rfloor} y^i q_{2i}(2\vartheta_y) A_0 = 0$, which satisfies (H) because $q_0$ is among the $q_{2i}$.
For $A_1$ apply the same argument to the odd part $f_o(x) = xA_1(x^2)$, using
$L x = x\,L|_{\vartheta \mapsto \vartheta+1}$.

Next, let $B(y) = \sum_{m=0}^{\infty} \beta_m y^m$ satisfy an equation with (H),
\[
  \sum_{j=0}^s y^j Q_j(\vartheta)B = 0, \qquad\text{that is,}\qquad
  \sum_{j=0}^{s} Q_j(m-j)\beta_{m-j} = 0 \quad (m \ge m_1),
\]
with $\deg Q_0 = d \ge \deg Q_j$. Let $\sigma \in \{0,\frac12\}$ and $g(u) = \sum_{m=0}^{\infty} \gamma_m u^m$ with
$\gamma_m = \beta_m/\Gamma(m+\sigma)$ (and $\gamma_0 = 0$ if $\sigma = 0$). Multiplying the recurrence by $1/\Gamma(m - s + \sigma)$ gives, for $m \ge m_1 + s$,
\[
  \sum_{j=0}^s \widetilde Q_j(m-j)\, \gamma_{m-j} = 0, \qquad
  \widetilde Q_j(m - j) = Q_j(m-j)\, (m - s + \sigma)_{s-j},
\]
with $\deg \widetilde Q_0 = d + s$ and $\deg \widetilde Q_j \le d + s - j$. Hence
$\sum_{j=0}^{s} u^j \widetilde Q_j(\vartheta) g$ is a polynomial, and in this operator the coefficient of the
highest power $\vartheta^{d+s}$ is the constant leading coefficient of $\widetilde Q_0$. Dividing by it
expresses $\vartheta^{d+s}g$ as a combination of $g, \dots, \vartheta^{d+s-1}g$ with polynomial
coefficients, plus a polynomial: this is \eqref{eq:Gsys} for $g$, with a companion matrix. Taking
$B = A_0$, $\sigma = 0$ gives $u h_0 = g$, and $\vartheta(u^{-1}g) = u^{-1}(\vartheta - 1)g$ transfers
\eqref{eq:Gsys} to $h_0$ (the inhomogeneous term becomes $u^{-1}b(u)$, a polynomial because $g$, hence
$b$, vanishes at $u = 0$); $B = A_1$, $\sigma = \frac12$ gives $h_1 = g$ directly.
\end{proof}

\subsection{The splitting}
Let
\[
  \theta_\chi(u) = \sum_{k \in \ZZ} \chi(k) k^\delta e^{-k^2u}, \qquad
  \hat\chi(n) = \sum_{a \bmod q} \chi(a) e^{2\pi i a n/q}.
\]
Poisson summation applied to $x^\delta e^{-ux^2}$ on each residue class modulo $q$ gives the theta transformation
\begin{equation}\label{eq:thetatrans}
  \theta_\chi(u) = c_\delta\, u^{-1/2-\delta} \sum_{n \in \ZZ} \hat\chi(n)\, n^\delta e^{-\pi^2 n^2/(q^2u)},
\end{equation}
where $c_\delta = (-i)^\delta \pi^{1/2+\delta} q^{-1-\delta}$.
Since $\hat\chi(-n) = (-1)^\delta\hat\chi(n)$, the terms $\pm n$ are equal.

\begin{proposition}\label{prop:split}
For every $u_0 > 0$,
\[
  S = \sum_{k=k_0}^{\infty} w(k)\, I_k + c_\delta \sum_{n=1}^{\infty} \hat\chi(n)\, n^\delta J_n + M,
\]
where
\begin{align*}
  I_k &= \int_{u_0}^\infty e^{-k^2u} h(u)\dd u, \\
  J_n &= \int_0^{u_0} u^{-1/2-\delta} e^{-\pi^2n^2/(q^2u)} h(u)\dd u, \\
  M &= \int_0^{u_0} \Bigl( [\delta = 0]\bigl(\tfrac12 c_0\hat\chi(0) u^{-1/2} - \tfrac12\chi(0)\bigr)
       - \sum_{k=1}^{k_0-1} w(k) e^{-k^2u} \Bigr) h(u)\dd u .
\end{align*}
\end{proposition}

\begin{proof}
By Lemma~\ref{lem:kernel}(2),
\[
  S = \sum_{k=k_0}^{\infty} w(k)\int_0^\infty e^{-k^2u}h(u)\dd u.
\]
Split each integral at $u_0$. On $[u_0,\infty)$ there is nothing more to do. On $(0, u_0]$ we may sum under the
integral sign: $\sum_{k\in\ZZ} |k|^\delta e^{-k^2u} = O(u^{-(1+\delta)/2})$ as $u \to 0$, and $h(u) = O(u^{\delta/2})$
(for $\delta = 1$ the coefficients start at $n = 3$), so the dominating function $O(u^{-1/2})$ is integrable. Finally
\[
  \sum_{k=k_0}^{\infty} w(k) e^{-k^2u} = \tfrac12\bigl(\theta_\chi(u) - [\delta=0]\chi(0)\bigr) - \sum_{k=1}^{k_0-1} w(k)e^{-k^2u},
\]
and \eqref{eq:thetatrans} turns $\frac12\theta_\chi$ into the $n = 0$ term (present only for $\delta = 0$)
plus $c_\delta u^{-1/2-\delta}\sum_{n=1}^{\infty}\hat\chi(n)n^\delta e^{-\pi^2n^2/(q^2u)}$.
\end{proof}

\subsubsection{Truncation} By Lemma~\ref{lem:kernel}(1), $|h(u)| \le C(1+u^{-1/2})(1+u)e^{R^2u}$, so
\[
  |I_k| \le C' \frac{e^{-(k^2 - R^2)u_0}}{k^2 - R^2}, \qquad
  |J_n| \le C'' e^{-\pi^2n^2/(q^2u_0)} \quad (n \ge 1),
\]
and the tails beyond the first neglected term are bounded by a geometric series. Hence $p$ bits need
\begin{equation}\label{eq:counts}
  k \le k_{\max} \approx \sqrt{\frac{p\log 2}{u_0}}, \qquad n \le n_{\max} \approx \frac{q}{\pi}\sqrt{p\log 2\cdot u_0},
\end{equation}
that is $O(\sqrt p)$ terms on each side. Note that $k_{\max}n_{\max} \approx q\,p\log 2/\pi$ does not
depend on $u_0$; the choice of $u_0$ only balances the two sides.

\subsection{The $k$-side}

\begin{lemma}\label{lem:kside}
Let $k \ge k_0$, $K = k^2$ and $x = Ku_0$. Then
\begin{equation}\label{eq:Ik}
\begin{split}
  I_k &= e^{-x} \sum_{m=1}^{\infty} a_{2m} \frac{e_m(x)}{K^m} + \erfc(\sqrt x)\, \frac{A_1(1/K)}{k} \\
      &\qquad + \frac{2\sqrt{u_0}}{\sqrt\pi}\, e^{-x} \sum_{m=1}^{\infty} \frac{a_{2m+1}}{K^m}
        \sum_{j=0}^{m-1} \frac{(2x)^j}{3\cdot5\cdots(2j+1)} .
\end{split}
\end{equation}
\end{lemma}

\begin{proof}
Termwise integration (Tonelli) gives $I_k = \sum_{n=2}^{\infty} a_n k^{-n}Q(n/2, x)$; insert $Q(m, x) = e^{-x}e_m(x)$ and
\[
  Q(m+\tfrac12, x) = \erfc\sqrt x + e^{-x}\sum_{j=0}^{m-1} \frac{x^{j+1/2}}{\Gamma(j+\frac32)},
\]
and use $\sqrt{x}/k = \sqrt{u_0}$.
\end{proof}

For rational $u_0$ the two series in \eqref{eq:Ik} have rational P-recursive terms: the sequences
$K^{-m}e_m(x)$ and $K^{-m}\sum_{j=0}^{m-1}(2x)^j/(3\cdots(2j+1))$ satisfy first-order inhomogeneous recurrences,
and $a_n$ is P-recursive. Their terms are bounded by $|a_n| K^{-n/2}$ times $e^{x}$ (which the prefactor
$e^{-x}$ cancels), so $O(p/\log(k/R))$ terms suffice, with $O(\log p)$-bit coefficients since $K, x = O(p)$.
$A_1(1/K)$ is a holonomic function at a rational point inside its disk of convergence, $\erfc\sqrt x$ is
needed only to absolute precision $2^{-p}$, and $e^{-x}$ is cheap. So each $I_k$ costs $O(\Mul(p)\log^2p)$.

\subsection{The dual side}\label{sec:dualgen}
For a function $g$ continuous on $[0, u_0]$, $t \in \QQ$ and $c > 0$ put
\[
  Z_t[g](c) = \int_0^{u_0} u^{-t} e^{-c/u} g(u)\dd u, \qquad E = E(c) = e^{-c/u_0}.
\]
If $g$ is $C^1$, then
\[
  \frac{\mathrm{d}}{\mathrm{d}u}\bigl(u^{1-t}e^{-c/u}g\bigr) = u^{-t}e^{-c/u}\Bigl((1-t)g + \frac cu g + \vartheta g\Bigr)
\]
and $u^{1-t}e^{-c/u} \to 0$ as $u \to 0^+$, so
\begin{equation}\label{eq:ibp}
  c\, Z_{t+1}[g] = E\, u_0^{1-t} g(u_0) - (1-t)\, Z_t[g] - Z_t[\vartheta g],
\end{equation}
and differentiation under the integral sign gives
\begin{equation}\label{eq:dc}
  -\partial_c Z_t[g] = Z_{t+1}[g].
\end{equation}
Also $Z_t[u^j g] = Z_{t-j}[g]$. Write $C_t = Z_t[1]$; by \eqref{eq:ibp},
$c\,C_{t+1} = E u_0^{1-t} - (1-t) C_t$.

\begin{lemma}\label{lem:moments}
Let $G$ satisfy \eqref{eq:Gsys}, $\vartheta G = A(u)G + b(u)$, with $A = \sum_{s=0}^{d} A_s u^s$ and
$b = \sum_{s=0}^{e} b_s u^s$. Fix $t$, and let $Y(c)$ be the vector formed by
$W_{t-s} = Z_{t-s}[G](c)$ for $0 \le s \le d$, by $C_{t-s}(c)$ for $0 \le s \le e$, by $E(c)$, and by the $r$ numbers
$E(c)\, (G(u_0))_i$. Then $Y$ satisfies a system \eqref{eq:fuchsys} whose matrices have entries in $\QQ(u_0^{1/2})$
depending only on $A$, $b$, $t$ and $u_0$, and $Y$ is analytic in $\operatorname{Re} c > 0$. (We take $u_0$ to be the
square of a rational number, so that the entries are rational; otherwise the bit-burst algorithm works over
$\QQ(u_0^{1/2})$ at a constant factor.)
\end{lemma}

\begin{proof}
By \eqref{eq:ibp} applied componentwise and \eqref{eq:Gsys},
\[
  c\, W_{t+1} = E\, u_0^{1-t} G(u_0) - (1-t) W_t - \sum_{s=0}^d A_s W_{t-s} - \sum_{s=0}^{e} b_s C_{t-s}.
\]
By \eqref{eq:dc}, $c\,\partial_c W_t = -c\,W_{t+1}$, which is the right-hand side above with the opposite sign,
a combination of components of $Y$ with constant coefficients; $c\,\partial_c W_{t-s} = -c\, W_{t-s+1}$ for
$s \ge 1$ is $c$ times a component of $Y$. The same holds for the $C_{t-s}$, using
$c\, C_{t+1} = E u_0^{1-t} - (1-t)C_t$, and for the last block, since $c\,\partial_c E = -(c/u_0)E$.
Analyticity: the integrals converge locally uniformly for $\operatorname{Re} c > 0$.
\end{proof}

By Proposition~\ref{prop:split} and \eqref{eq:kernels},
\[
  J_n = Z_{1/2+\delta}[h_0](c) + Z_{1+\delta}[h_1](c), \qquad c = \frac{\pi^2n^2}{q^2},
\]
so by Lemmas~\ref{lem:kernel}(3) and~\ref{lem:moments} the dual term $J_n$ is a
fixed linear form in $Y(c)$ for one system \eqref{eq:fuchsys} (two blocks, for $h_0$ and $h_1$). We evaluate
it by the bit-burst algorithm from a rational point $c_0$ with $|c - c_0| \le 2^{-8}c$, say.

\subsubsection{Starting values} Let
\[
  D_s(c) = \int_0^{u_0}u^{s-1}e^{-c/u}\dd u = c^s\,\Gamma(-s, c/u_0).
\]
Then $Z_t[g](c_0) = \sum_{j=0}^{\infty} g_j D_{j-t+1}(c_0)$ for $g = \sum_{j=0}^{\infty} g_j u^j$, and \eqref{eq:ibp} with $g = u^s$ gives
\begin{equation}\label{eq:Drec}
  s\, D_s + c\, D_{s-1} = E\, u_0^s .
\end{equation}
For $t \in \frac12 + \ZZ$ the exponents $s = j - t + 1$ are half-integers, and \eqref{eq:Drec} expresses every
$D_s$ as $\alpha_s D_{-1/2} + \beta_s E$ with $D_{-1/2} = c^{-1/2}\Gamma(\frac12, c/u_0)$; for $t \in \ZZ$ they are integers,
and every $D_s$ with $s \ge 0$ is $\alpha_sE_1(c/u_0) + \beta_sE$, while $D_{-1} = E/c$ and the $D_s$ with $s \le -2$
are rational multiples of $E$. The coefficients
$\alpha_s, \beta_s$ satisfy first-order recurrences with coefficients in $\QQ(c_0, u_0^{1/2})$. Therefore every
component of $Y(c_0)$ is a combination of $\Gamma(\frac12, c_0/u_0)$ or $E_1(c_0/u_0)$, $E(c_0)$ and the constants
$G(u_0)$, with coefficients given by two series with rational P-recursive terms. Individually, each
$D_s = \alpha_sD_* + \beta_sE$ (with $D_*$ the base value) cancels by up to $e^{c_0/u_0}$. In the weighted sums this
does not accumulate: since $|g_j| \le CR^{2j}/j!$ and $|\alpha_s| \le c_0^{s+O(1)}/\Gamma(s+1)$, the terms of
$\sum_{j=0}^{\infty} g_j\alpha_{j-t+1}$ are of size
\[
  \frac{(R^2c_0)^j}{(j!)^2} \le e^{2R\sqrt{c_0}},
\]
and $E\beta_s = D_s - \alpha_s D_*$ gives the same for the $\beta$-series, while the result is of size $E$ times a factor
polynomial in $c_0$. So $O(\sqrt{c_0} + \log p) = O(\sqrt p)$ guard bits suffice, and the cost is $O(\Mul(p)\log^2p)$.
The base values $\Gamma(\frac12, y)$ and $E_1(y) = -\gamma - \log y + \Ein(y)$ at the rational point $y = c_0/u_0 = O(p)$
cost $O(\Mul(p)\log^2 p)$ with $O(p)$ guard bits for the alternating series.

\subsubsection{Cost of a dual term} The system has fixed dimension, its matrices depend only on $t$, $A$, $b$ and
$u_0$ (not on $n$), and $c_0$ has $O(1)$ significant bits and $O(\log p)$ bits in all. We take $c_0 \le c$, so that the
factor $E(c)/E(c_0) \le 1$ does not amplify errors in the starting values; the other solutions of
\eqref{eq:fuchsys} grow at most like $e^{O(\sqrt c)}$ along the path, which costs $O(\sqrt p)$ guard bits. The
bit-burst steps from $c_0$ to $c = \pi^2n^2/q^2$ cost $O(\Mul(p)\log^2p)$, and $\pi^2$ is computed once.
The a priori bound needed for the Taylor tails holds since $|Y(c)| \le C e^{O(1)}$ on
$|c - c_0| \le \frac34 c_0$. So each $J_n$ costs $O(\Mul(p)\log^2p)$.

\subsection{Main term and constants}
$M$ is a finite combination of the integrals
\[
  \int_0^{u_0} u^{-1/2} h\dd u, \qquad \int_0^{u_0} h\dd u, \qquad \int_0^{u_0}e^{-k^2u}h\dd u \quad (k < k_0);
\]
expanding the integrands in powers of $u^{1/2}$ gives series in $u_0^{1/2}$ with rational P-recursive terms that
converge like $e^{O(u_0)}$-type entire series. Together with $\pi$, $\gamma$ (only if $E_1$ occurs), the values
$\vartheta^iG_\varepsilon(u_0)$, and the $O(\sqrt p)$ exponentials $e^{-k^2u_0}$ and $e^{-c/u_0}$ (each
$O(\Mul(p)\log^2 p)$, or by multiplicative recurrences), all of this costs $O(\sqrt p\,\Mul(p)\log^2p)$.

\begin{proof}[Proof of Theorem~\ref{thm:main}]
By Proposition~\ref{prop:split} and \eqref{eq:counts}, $S$ is a combination of $O(\sqrt p)$ terms $I_k$
and $J_n$ plus $M$, each computable to $p + O(\log p)$ bits in $O(\Mul(p)\log^2p)$; the coefficients
$w(k)$ and $c_\delta\hat\chi(n)n^\delta$ lie in $\QQ(\zeta_{\operatorname{lcm}(q,4)}, \pi^{1/2})$ and are bounded by $O(\sqrt p)$.
Arithmetic in a fixed cyclotomic field costs $O(1)$ multiplications of rationals per operation.
\end{proof}

\begin{algorithm}[Summary]\label{alg:main}
Input $f$ (by its operator and initial coefficients), $\chi$, $k_0$, $p$.
\begin{enumerate}
\item Split $\chi$ by parity; for each part, form $h_0, h_1$ and the systems \eqref{eq:Gsys}; choose $u_0 \in \QQ$.
\item For $k_0 \le k \le k_{\max}$, evaluate $I_k$ by \eqref{eq:Ik} with binary splitting.
\item For $1 \le n \le n_{\max}$, with $c = \pi^2n^2/q^2$: compute $Y(c_0)$ at a rational $c_0 \approx c$ from
the starting value series; bit-burst to $c$; $J_n$ is a linear form in $Y(c)$.
\item Compute $M$; combine as in Proposition~\ref{prop:split}.
\end{enumerate}
\end{algorithm}

\begin{remark}[Practical points]\label{rem:practical}
(1) The terms are independent, so the algorithm parallelizes trivially, with $O(p)$ memory per term. (2) The
$k$-side terms need only relative precision $p - (k^2u_0)\log_2 e$ and the dual terms $p - (c/u_0)\log_2e$,
which also shortens the series. (3) For $u_0 = 4^{-a}$ the arguments $\sqrt x = k\sqrt{u_0}$ and $c_0/u_0$ are
short dyadic numbers, so $\erfc\sqrt x$ and $\Gamma(\frac12, c_0/u_0)$ are cheap binary splitting series. (4) A
dual term is typically several dozen times as expensive as a $k$-term, which pushes the optimal $u_0$ below
$1$ (in our applications $u_0 = \frac1{16}$).
\end{remark}

\begin{remark}[Extensions]\label{rem:ext}
(1) Algebraic coefficients $a_n$ in a fixed number field work the same way. (2) The terms $k < k_0$ can be
added separately when $f$ is analytic at the points $1/k$ (analytic continuation along a path avoiding the
singularities, $O(\Mul(p)\log^2p)$ per point), so that $k_0 = 1$ is allowed. (3) A linear term $a_1 x$ contributes
$a_1\sum_{k=k_0}^{\infty}\chi(k)/k$, which converges when $\sum_{a=1}^{q}\chi(a) = 0$ and is a combination of logarithms and
$\pi$. (4) The weights $e^{-k^2u}$ can be replaced by $e^{-Q(k)u}$ for a positive definite quadratic form in
several variables, which gives the corresponding statement for lattice sums $\sum_{k \in \ZZ^d} f(1/Q(k))$
(with Ewald's theta transformation), at the price of $O(p^{d/2})$ terms. (5) Weights $\chi(k) = e^{2\pi i\alpha k}$
with real algebraic $\alpha$ need not be periodic, but the reduction to even weights still applies:
$\chi_0(k) = \cos 2\pi\alpha k$, and $\chi_1(k)k = i k\sin 2\pi\alpha k$ is even. Poisson summation gives
\[
  \sum_{k\in\ZZ}\cos(2\pi\alpha k)\,e^{-k^2u} = \sqrt{\frac\pi u}\,\sum_{n\in\ZZ}e^{-\pi^2(n-\alpha)^2/u},
\]
and its derivative in $\alpha$ for the odd part, so the dual terms are the same moments at the points $c = \pi^2(n-\alpha)^2$, $n \in \ZZ$, and
there is no term with $c = 0$ unless $\alpha \in \ZZ$. This covers sums of Lerch type $\sum_{k=k_0}^{\infty} e^{2\pi i\alpha k}f(1/k)$.
\end{remark}

\begin{remark}[Uniformity]\label{rem:uniform}
For $f(x) = x^s$ the kernel is a monomial and every term is a single incomplete gamma function; this is the
approximate functional equation for $L(s,\chi)$, and \cite{Joh21} shows that the cost remains
$p^{3/2+o(1)}$ even when $s$ grows with $p$, which gives exact Bernoulli numbers $B_n$ in $n^{3/2+o(1)}$.
Theorem~\ref{thm:main} is uniform only for fixed $f$: its constants grow with the order and degree of the
equation for $f$ and with the size of its coefficients. We have not found a family of exactly computable
integer sequences that would profit from Theorem~\ref{thm:main} in the same way.
\end{remark}

\subsection{Logarithmic weights}

\begin{theorem}\label{thm:log}
Under the hypotheses of Theorem~\ref{thm:main}, for every fixed integer $j \ge 1$,
\[
  S^{(j)} = \sum_{k=k_0}^{\infty}\chi(k)\,(\log k)^j f(1/k)
\]
can be computed to $p$ bits in $O(p^{1/2}\Mul(p)\log^2p)$ bit operations.
\end{theorem}

\begin{proof}[Proof sketch]
Let
\[
  S(\sigma) = \sum_{k=k_0}^{\infty} \chi(k) k^{-2\sigma} f(1/k), \qquad\text{so that}\qquad
  S^{(j)} = \Bigl(-\frac12\partial_\sigma\Bigr)^jS(\sigma)\Big|_{\sigma=0}.
\]
Everything above holds for complex $\sigma$ near $0$ with the kernel
\[
  h_\sigma(u) = \sum_{n=2}^{\infty} a_n \frac{u^{n/2+\sigma-1}}{\Gamma(n/2+\sigma)},
\]
and all pieces of Proposition~\ref{prop:split}
are analytic in $\sigma$, so $S^{(j)}$ is the sum of their $\sigma$-derivatives. We indicate why each derivative
is computable within the same bound.

\emph{Kernel.} We have
\[
  \partial_\sigma^i h_\sigma = u^\sigma\sum_{l=0}^{i}\binom il (\log u)^{i-l}\, \partial_\sigma^l g_\sigma, \qquad
  g_\sigma = \sum_{n=2}^{\infty} a_n \frac{u^{n/2-1}}{\Gamma(n/2+\sigma)}.
\]
In the proof of
Lemma~\ref{lem:kernel}(3) the Pochhammer factors $(m - s + \sigma_0 + \sigma)_{s-j}$ are polynomials in $\sigma$, so
$\vartheta G_\sigma = A_\sigma(u)G_\sigma + b_\sigma(u)$ with $A_\sigma$ polynomial in $\sigma$; the inhomogeneous polynomial $b_\sigma$
comes from the first few coefficients $\beta_m/\Gamma(m+\sigma_0+\sigma)$ and is entire in $\sigma$, its $\sigma$-jets being
rational combinations of polygamma values at integers and half-integers (that is, of rationals, $\gamma$, $\log2$ and
$\zeta(2), \dots, \zeta(j)$). Differentiating in $\sigma$ shows that the jets $(\partial_\sigma^lG_\sigma|_{\sigma=0})_{l\le j}$
satisfy a block-triangular system of the form \eqref{eq:Gsys} whose inhomogeneous part has such constant coefficients.

\emph{Dual side.} The moments $Z^{(i)}_t[g] = \int_0^{u_0}u^{-t}(\log u)^ie^{-c/u}g\dd u$ satisfy
\[
  c Z^{(i)}_{t+1}[g] = E u_0^{1-t}(\log u_0)^i g(u_0) - (1-t)Z^{(i)}_t[g] - Z^{(i)}_t[\vartheta g] - iZ^{(i-1)}_t[g]
\]
and $-\partial_c Z^{(i)}_t = Z^{(i)}_{t+1}$, so Lemma~\ref{lem:moments} holds for the enlarged, block-triangular
system (with the forcing constants $E u_0^{1-t}(\log u_0)^i G(u_0)$ as further state components). Starting values need
$\partial_s^iD_s$; differentiating \eqref{eq:Drec} reduces them to $\partial_s^i\Gamma(s, y)$ at $s \in \{0, \pm\frac12\}$
(the values at $s = -1$ follow from the differentiated recurrence at $s = 0$) and rational $y$; these are series in $y$ with rational
P-recursive terms plus the constants $\gamma$, $\zeta(2), \dots, \zeta(i+1)$, $\log 2$, $\log y$, all computable in
$O(\Mul(p)\log^2 p)$.

\emph{$k$-side.} $\partial_\sigma^i[k^{-n-2\sigma}Q(n/2+\sigma, x)]$ involves $(\log k)^l$ (one logarithm per $k$, $O(\sqrt p)$ in all) and
$\partial_s^l Q(s, x)$ at $s = n/2$. Differentiating $Q(s+1, x) = Q(s, x) + x^se^{-x}/\Gamma(s+1)$ gives
P-recursive inhomogeneous recurrences in $s$ (with polygamma values $\psi^{(l)}(m)$, $\psi^{(l)}(m+\frac12)$, which
are rational plus constants), with initial values computable as above.

\emph{Main term.} Integrals $\int_0^{u_0}u^{a}(\log u)^i\dd u$ are explicit.
\end{proof}

\begin{example}
$\sum_{k=2}^{\infty}\frac{\log k}{k^2-1}$ is $S^{(1)}$ for the rational function $f(x) = x^2/(1-x^2)$; it is also an
instance of Theorem~\ref{thm:main} without logarithms (Section~\ref{sec:AG}). Likewise
$\sum_{k=1}^{\infty} \frac{\log k}{k(k+1)}$ has $f(x) = x^2/(1+x)$. Derivatives $L'(s,\chi)$ at integers
(and the Stieltjes constants) are covered by \cite{Joh21}.
\end{example}

\subsection{Functions regular at infinity}\label{sec:regular}

Let $L = \sum_{i=0}^r p_i(x)\,\partial_x^i$ with $p_i \in \Qbar[x]$, $p_r \ne 0$, and suppose that $x = \infty$ is an
ordinary point or a regular singular point of $L$, i.e.\ that $t = 1/x$ is such a point of the transformed
operator. Let $R$ exceed the absolute values of the finite singular points of $L$. By Frobenius's method
\cite[\S16]{Inc26}, the solutions of $L$ on $(R, \infty)$ have a basis of functions
$x^{-\sigma}\sum_{j=0}^{J}(\log x)^j\varphi_j(1/x)$, where $\sigma$ is a root of the indicial polynomial at infinity and
the $\varphi_j$ are power series with coefficients in a number field, convergent for $|t| < 1/R$. A solution $F$ is
\emph{defined over $\Qbar$} if its coordinates in such a basis are algebraic (this does not depend on the choice
of the basis among those with algebraic coefficients). Collecting terms gives \eqref{eq:frob}, with exponents
$\sigma \in \Sigma$ pairwise incongruent modulo $\ZZ$ and power series $f_{\sigma,j}$ with coefficients in a number field.

\begin{theorem}\label{thm:general}
Let $L$ and $R$ be as above, let $F$ be a solution of $L$ on $(R,\infty)$ defined over $\Qbar$ with
$F(x) = O(x^{-1-\varepsilon})$ as $x \to +\infty$ for some $\varepsilon > 0$, let $\chi \colon \ZZ \to \QQ(\zeta_q)$
have period $q$, and let $k_0 > R$ be an integer. Then $\sum_{k=k_0}^{\infty}\chi(k)F(k)$ can be computed to $p$ bits
in $O(p^{1/2}\Mul(p)\log^2 p)$ bit operations.
\end{theorem}

\begin{proof}[Proof sketch]
\emph{The pieces.} Write $B_\sigma = x^{-\sigma}\sum_{j=0}^{J}(\log x)^jf_{\sigma,j}(1/x)$ for the terms of \eqref{eq:frob}. Each
$B_\sigma$ is obtained from $F$ by a polynomial in the monodromy operator $M$ around $t = 0$ (the projection onto the
generalized eigenspace for the eigenvalue $e^{2\pi i\sigma}$), so it is a solution of $L$. Since
\[
  e^{-2\pi i\sigma}MB_\sigma = x^{-\sigma}\sum_{j=0}^{J}(\log x - 2\pi i)^jf_{\sigma,j}(1/x),
\]
the functions $x^{-\sigma}\sum_{j=0}^{J}(\log x - 2\pi i m)^jf_{\sigma,j}(1/x)$, $0 \le m \le J$, are solutions of $L$ too, and solving for the
coefficients of the powers of $m$ expresses each $f_{\sigma,j}(t)$ as a polynomial in $\log t$ whose coefficients are
of the form $t^{-\sigma}y(1/t)$ with $y$ a solution of $L$. Moderate growth at $t = 0$ is preserved by sums and
products, so $f_{\sigma,j}$ is a holonomic power series annihilated by an operator with a regular singular point at
$t = 0$, which satisfies (H) by Fuchs's criterion. Replacing $\sigma$ by $\sigma + m$ and $f_{\sigma,j}$ by
$t^{-m}f_{\sigma,j}$ we may assume that the exponents are distinct and that for each $\sigma$ some $f_{\sigma,j}$
has a nonzero constant term. The functions $x^{-\sigma}(\log x)^j$ with distinct pairs $(\sigma, j)$ do not cancel
asymptotically (for fixed real part, the factors $x^{-i\operatorname{Im}\sigma}$ are linearly independent almost
periodic functions of $\log x$), so the decay condition gives $\operatorname{Re}\sigma > 1$ for every $\sigma$ that occurs.
It therefore suffices to compute
\begin{align*}
  S_{\sigma,j} &= \sum_{k=k_0}^{\infty}\chi(k)\,k^{-\sigma}(\log k)^jf(1/k) = (-\partial_\tau)^j\, S_\sigma(\tau)\big|_{\tau=0},\\
  S_\sigma(\tau) &= \sum_{k=k_0}^{\infty}\chi(k)\,k^{-\sigma-\tau}f(1/k),
\end{align*}
for an algebraic $\sigma$ with $\operatorname{Re}\sigma > 1$ and a power series $f = \sum_{n=0}^{\infty}a_nx^n$ satisfying (H)
with radius of convergence $> 1/k_0$.

\emph{The splitting with a shifted exponent.} For $\sigma$ near $0$, the proof of Theorem~\ref{thm:log} uses
the kernel $\sum_{n=2}^{\infty}a_nu^{n/2+\sigma-1}/\Gamma(n/2+\sigma)$. It goes through for any fixed algebraic $\sigma$
with $\operatorname{Re}\sigma > 1$, with $(\sigma+\tau)/2$ in place of that $\sigma$ (and $\tau$ near $0$), and with the
following changes. (i) The Laplace representation
\[
  k^{-n-\sigma} = \frac{1}{\Gamma(\frac{n+\sigma}2)}\int_0^\infty u^{(n+\sigma)/2-1}e^{-k^2u}\dd u
\]
holds since $\operatorname{Re}(n + \sigma) > 0$, and the integrals in Proposition~\ref{prop:split} converge at $u = 0$
because the dominating function is now of order $u^{(\operatorname{Re}\sigma-3)/2}$. (ii) Lemma~\ref{lem:kernel}(3) holds with
$\Gamma(m+\sigma_0)$ for any $\sigma_0$ in place of $\sigma_0 \in \{0, \frac12\}$; the kernel is $u^{\sigma/2}$ times
two entire holonomic functions (one for each parity of $n$), with the constants $1/\Gamma(\sigma/2)$ and
$1/\Gamma(\frac{\sigma+1}2)$ factored out. (iii) On the dual side the moment indices $t$ lie in $-\frac\sigma2 + \frac12\ZZ$, and
\eqref{eq:Drec} reduces all starting values to $E$ and one incomplete gamma value $\Gamma(a, c_0/u_0)$ per class
of $a$ modulo $\ZZ$ (and its derivatives in $a$ for $j > 0$), with $a$ algebraic; on the $k$-side the terms involve
$k^{-\sigma}$ and the regularized incomplete gamma functions $Q(\frac{n+\sigma}{2}, k^2u_0)$, which satisfy the same
recurrence in $n$ as before. (iv) The main term consists of explicit integrals $\int_0^{u_0}u^{b}(\log u)^i\dd u$ with
algebraic $b$.

\emph{Cost.} All recurrences now have coefficients in a fixed number field containing $\sigma$, which costs a
constant factor. The new constants are $\Gamma$ and its derivatives at the points $\sigma/2$ and $(\sigma+1)/2$,
and the new per-term quantities are $k^{-\sigma} = e^{-\sigma\log k}$ and $\Gamma(a, y)$ for algebraic $a$ and
rational $y = O(p)$. All of these cost $O(\Mul(p)\log^2p)$: $\log k$ and the exponential by binary splitting and the
bit-burst algorithm; $\Gamma(a, y)$, as in the rational case, by binary splitting of the lower series
$\gamma(a, y) = y^ae^{-y}\sum_{i=0}^{\infty}y^i/(a)_{i+1}$, with $O(p)$ guard bits; and $\Gamma(a)$ as
$\gamma(a, N) + \Gamma(a, N)$ with $N = O(p)$, both by binary splitting over the number field. There are
$O(\sqrt p)$ such per-term quantities, so the bound of Theorem~\ref{thm:main} is unchanged.
\end{proof}

\begin{remark}
For $f = 1$, $\chi$ a Dirichlet character and $j = 0$, Theorem~\ref{thm:general} is the statement of \cite{Joh21} for
$L(\sigma, \chi)$ at a fixed algebraic point. Conversely, Theorem~\ref{thm:general} with $\chi = 1$ already implies the
version with periodic $\chi$ and with additional factors $(\log k)^j$, by the closure properties mentioned in the
introduction: $F(qx+a)$ and $(\log x)^jF(x)$ are again holonomic, defined over $\Qbar$ and regular at infinity.
\end{remark}

\begin{remark}[Integrals against Bose--Einstein weights]\label{rem:bose}
Integrals of holonomic functions over intervals with algebraic endpoints can be computed in quasilinear time,
but integrals against weights that are not holonomic are not covered by this. Some of them are sums in our
class. Let $G(t) = \sum_{n=1}^{\infty}g_nt^n/n!$ be an entire holonomic function of exponential type $\tau$ with algebraic
coefficients. Its Laplace transform $\int_0^\infty G(t)e^{-kt}\dd t = \sum_{n=1}^{\infty}g_nk^{-n-1}$ is then holonomic,
defined over $\Qbar$ and analytic at infinity for $k > \tau$, and
\[
  \int_0^\infty \frac{G(t)}{e^t - 1}\dd t = \sum_{k=1}^{\infty}\int_0^\infty G(t)e^{-kt}\dd t,
\]
when the integral converges (the finitely many terms with $k \le \tau$ are integrals of holonomic functions), with
alternating signs for the Fermi--Dirac weight $1/(e^t+1)$. So Planck-type integrals
of such $G$ can be computed in time $\widetilde O(p^{3/2})$; for example
\[
  \int_0^\infty \frac{t\,J_0(at)}{e^t-1}\dd t = \sum_{k=1}^{\infty}\frac{k}{(k^2+a^2)^{3/2}}.
\]
Conversely, the kernel $h$ of Section~\ref{sec:general} shows that every sum \eqref{eq:S} is an integral of a
holonomic function (entire up to a factor $u^{-1/2}$) against a theta series, and the first-order Euler--Maclaurin
formula
\[
  \sum_{k=1}^{\infty}F(k) = \int_1^\infty F(x)\dd x + \frac12F(1) + \int_1^\infty\Bigl(\{x\}-\frac12\Bigr)F'(x)\dd x
\]
does the same with the sawtooth weight. Gautschi and Milovanovi\'c \cite{GM85} use the opposite direction at fixed precision, summing series by Gaussian
quadrature with Bose--Einstein and Fermi--Dirac weights.
\end{remark}

%======================================================================
\section{Khinchin's constant}\label{sec:khinchin}

\subsection{A dilogarithm series}

\begin{lemma}[{Bailey, Borwein and Crandall \cite[\S3]{BBC97}}]\label{lem:dilog}
\[
  \log K \log 2 = \log^2 2 + \Li_2(-\tfrac12) + \tfrac12\sum_{k=2}^{\infty}(-1)^k \Li_2\Bigl(\frac{4}{k^2}\Bigr).
\]
\end{lemma}

This is the first formula of Section~3 of \cite{BBC97}, where it is attributed to Zagier and derived from an
integral representation of $\log K$ via the Euler product for $\sin z/z$; it can also be derived directly from
\eqref{eq:logK} by telescoping. Its advantage here is that the coefficients $4^n/n^2$ are hypergeometric, whereas
those of $f(x) = -\log(1-x)\log(1+x)$, the power series of \eqref{eq:logK} (Section~\ref{sec:logform}), contain
alternating harmonic numbers.

\subsection{The explicit formula}

For $k \ge 3$,
\[
  \Li_2\Bigl(\frac4{k^2}\Bigr) = \int_0^\infty e^{-k^2u} G(u)\dd u, \qquad
  G(u) = \frac{\Eip(4u)}{u} = \sum_{n=1}^{\infty}\frac{4^nu^{n-1}}{n\cdot n!},
\]
where $\vartheta\Eip(4u) = e^{4u} - 1$.
The alternating sign makes the relevant theta function
\[
  \theta_4(u) = \sum_{k\in\ZZ}(-1)^ke^{-k^2u} = 2\sqrt{\frac{\pi}{u}}\sum_{m=0}^{\infty}e^{-\pi^2(m+\frac12)^2/u},
\]
whose transform has no constant term (this is
\eqref{eq:thetatrans} with $\chi(k) = (-1)^k$, $q = 2$ and $\hat\chi(n) = 1 - (-1)^n$).

\begin{theorem}\label{thm:khinchin}
For every $u_0 > 0$,
\[
  \log K\log 2 = \log^2 2 + \Li_2(-\tfrac12) + \frac{\pi^2}{12} + \frac12\Bigl(M + \sum_{k=3}^{\infty}(-1)^kI_k + \sum_{m=0}^{\infty}J_m\Bigr),
\]
where, with $x_k = k^2u_0$, $c_m = \pi^2(m+\frac12)^2$ and $y_m = c_m/u_0$,
\begin{align*}
  I_k &= e^{-x_k}\sum_{n=1}^{\infty}\Bigl(\frac{4}{k^2}\Bigr)^n\frac{e_n(x_k)}{n^2}, \\
  J_m &= \frac{2\pi}{\sqrt{c_m}}\erfc(\sqrt{y_m})\bigl(\Ci(4\sqrt{c_m}) - \gamma - \log(4\sqrt{c_m})\bigr) \\
      &\qquad + \sqrt\pi\, e^{-y_m}\mathcal B(c_m), \\
  \mathcal B(c) &= \sum_{n=1}^{\infty}\frac{4^n}{n\cdot n!}\,\beta_n(c), \\
  \beta_n(c) &= \sum_{i=1}^{n}(-c)^{n-i}u_0^{i-\frac12}\frac{\Gamma(i-\frac12)}{\Gamma(n+\frac12)}, \\
  M &= \sum_{n=1}^{\infty}\frac{4^n}{n^2}\Bigl(\pi_n(u_0) - 4^{-n}\pi_n(4u_0)\Bigr) - \frac12\sum_{n=1}^{\infty}\frac{(4u_0)^n}{n^2\,n!},
\end{align*}
with $\pi_n(y) = 1 - e^{-y}e_n(y)$.
\end{theorem}

\begin{proof}
As in Proposition~\ref{prop:split}, with
\[
  \sum_{k=3}^{\infty}(-1)^ke^{-k^2u} = \tfrac12\theta_4(u) - \tfrac12 + e^{-u} - e^{-4u},
\]
we get
\[
  T_3 := \sum_{k=3}^{\infty}(-1)^k\Li_2(4/k^2) = \sum_{k=3}^{\infty}(-1)^k I_k + \sum_{m=0}^{\infty} J_m + M,
\]
where
\begin{gather*}
  I_k = \int_{u_0}^\infty e^{-k^2u}G\dd u, \qquad J_m = \sqrt\pi\int_0^{u_0}u^{-1/2}e^{-c_m/u}G\dd u, \\
  M = \int_0^{u_0}\Bigl(e^{-u} - e^{-4u} - \frac12\Bigr)G\dd u.
\end{gather*}
(Interchanges: on $[u_0,\infty)$ the series is dominated by
$e^{-5u}/u_0$; on $(0,u_0]$, $\sum_{k=0}^{\infty} e^{-k^2u} \le 1 + \sqrt{\pi/u}$ and $G$ is bounded.) Then
by Lemma~\ref{lem:dilog},
\[
  \log K\log2 = \log^22 + \Li_2(-\tfrac12) + \frac12\Bigl(\frac{\pi^2}6 + T_3\Bigr).
\]

$I_k$: termwise, $\int_{u_0}^\infty e^{-Ku}u^{n-1}\dd u = K^{-n}(n-1)!\,e^{-x}e_n(x)$.

$J_m$: termwise,
\[
  J_m = \sqrt\pi\sum_{n=1}^{\infty} \frac{4^n}{n\cdot n!}D_{n-1/2}(c)
\]
with $D_s$ as in Section~\ref{sec:dualgen}. By \eqref{eq:Drec},
\[
  D_{n-1/2} = \alpha_n D_{-1/2} + \beta_n E, \qquad \alpha_n = \frac{(-c)^n}{(\frac12)_n},
\]
with the $\beta_n$ as stated, and $D_{-1/2} = \sqrt{\pi/c}\erfc\sqrt{c/u_0}$ (substitute $u = c/s^2$). Finally
\[
  \sum_{n=1}^{\infty} \frac{4^n}{n\cdot n!}\alpha_n = \sum_{n=1}^{\infty}\frac{(-16c)^n}{n\,(2n)!} = 2\bigl(\Ci(4\sqrt c) - \gamma - \log(4\sqrt c)\bigr).
\]

$M$: termwise, with $\int_0^{u_0}e^{-au}u^{n-1}\dd u = a^{-n}(n-1)!\,\pi_n(au_0)$.
\end{proof}

All three series converge for every $u_0 > 0$: the $I_k$ series like $(4/k^2)^n$, and those for
$\mathcal B$ and $M$ like $(16c)^n/(2n)!$ and $(4u_0)^n/n!$. The formula was checked numerically to
$40$ digits for $u_0 = \frac14$ with independent code, and the implementation below agrees with the zeta
series method to $10^6$ digits. For computation we use
\[
  \Li_2(-\tfrac12) = -\Li_2(\tfrac13) - \tfrac12\log^2\tfrac32, \qquad
  \Li_2(\tfrac13) = \tfrac16\Li_2(\tfrac19) + \frac{\pi^2}{18} - \tfrac16\log^23.
\]

\subsection{Evaluation}\label{sec:kheval}
We evaluate the three parts of Theorem~\ref{thm:khinchin} as follows.

\subsubsection{$k$-side} With $K = k^2$, $P_n = (4/K)^ne_n(x)$ and $Q_n = (4u_0)^n/n!$,
\[
  P_{n+1} = \tfrac4K(P_n + Q_n), \qquad Q_{n+1} = \tfrac{4u_0}{n+1}Q_n, \qquad I_k = e^{-x}\sum_{n=1}^{\infty}P_n/n^2,
\]
a $2\times2$ triangular recurrence. Carrying $U_n = P_n/n$ instead, the recurrence
\[
  U_{n+1} = \frac{4}{K}\,\frac{nU_n + Q_n}{n+1}
\]
shares the denominator $K(n+1)$ with $Q$, and the sum $\sum_{n=1}^{\infty} U_n/n$ has weight $1/n$ only. For $u_0 = 4^{-a}$ the
leaves have about $\log_2 K + 2\log_2 n$ bits.

\subsubsection{Dual side} With $\nu_t = Z_t[\Eip(4u)]$, $Y_t = Z_t[e^{4u}]$ and $C_t = Z_t[1]$ (Section~\ref{sec:dualgen}),
\eqref{eq:ibp} gives
\begin{align*}
  c\,\nu_{t+1} &= E\,\Eip(4u_0)\,u_0^{1-t} - (1-t)\nu_t - Y_t + C_t, \\
  c\,Y_{t+1} &= E\,e^{4u_0}u_0^{1-t} - (1-t)Y_t - 4Y_{t-1}, \\
  c\,C_{t+1} &= E\,u_0^{1-t} - (1-t)C_t,
\end{align*}
and $J_m = \sqrt\pi\,\nu_{3/2}(c_m)$. So $J_m$ is a linear form in the solution of a system of order $4$
(state $\nu_t, Y_t, Y_{t-1}, C_t$; three forcing constants $\Eip(4u_0)$, $e^{4u_0}$, $1$) with coefficients
linear in $c$ and $t$; in the scaled variables $c^jZ_{t_0+j}$ the leading coefficient of the Taylor recurrence is $1$.
Starting values at a $64$-bit rational $c_0$ come from the series of Theorem~\ref{thm:khinchin},
\begin{gather*}
  \nu_{3/2} = \sum_{n=1}^{\infty}\frac{4^n}{n\cdot n!}D_{n-1/2}, \qquad Y_{3/2} = \sum_{n=0}^{\infty}\frac{4^n}{n!}D_{n-1/2}, \\
  Y_{1/2} = \sum_{n=0}^{\infty}\frac{4^n}{n!}D_{n+1/2}, \qquad C_{3/2} = D_{-1/2},
\end{gather*}
summed by one binary splitting with a three-term state ($\alpha_n$, $\beta_n$, and the factor $u_0^{n-1/2}$) and
seven sums. The bit-burst uses rational points $c_j$ with $64\cdot4^j$ fractional bits; the Taylor stage from
$c_{j-1}$ to $c_j$ sums
\[
  Z_t(c_{j-1}+\delta) = \sum_{l=0}^{\infty}\frac{(-\delta)^l}{l!}Z_{t+l}(c_{j-1}),
\]
with $Z_{t+l+1}$ obtained from the recurrences.

\subsubsection{Main term and constants} $M$, $\Eip(4u_0)$, $e^{4u_0}$ and $\Li_2(\frac19)$ are rational series;
$e^{-x_k}$ is propagated by $e^{-x_{k+1}} = e^{-x_k}e^{-(2k+1)u_0}$.

\subsection{Error bounds and convergence}\label{sec:khbounds}
All computations use ball arithmetic; the following bounds (with $U = u_0$, $E = e^{-c/U}$) make the result rigorous.

\begin{itemize}
\item \emph{Omitted terms.} Since $G(u) \le e^{4u}/U$ for $u \ge U$,
\[
  0 < I_k \le \frac{e^{-(k^2-4)U}}{(k^2-4)U}, \qquad 0 < J_m \le 2\sqrt{\pi U}\,G(U)\,e^{-\pi^2(m+\frac12)^2/U},
\]
and both tails are at most twice their first term. Thus $p$ bits need
\[
  k_{\max} \approx \sqrt{4 + p\log2/U}, \qquad m_{\max} \approx \frac{\sqrt{pU\log 2}}{\pi}:
\]
at $10^6$ digits and $U = \frac1{16}$, $k_{\max} = 6\,069$ and $121$ dual terms.
\item \emph{$k$-series.} The $n$-th term of $e^{-x}\sum_{n=1}^{\infty} P_n/n^2$ is
\[
  \Bigl(\frac4K\Bigr)^n\frac{\Pr[\mathrm{Poisson}(x) \le n-1]}{n^2},
\]
bounded by
\[
  \tau_n = \Bigl(\frac4K\Bigr)^n\min\Bigl(1, \frac{n\,x^{n-1}e^{-x}}{(n-1)!}\Bigr)\frac{1}{n^2}
\]
for $n - 1 < x$ and by $(4/K)^n/n^2$ beyond. For $U \le \frac14$ the ratios are at most $\frac12$ and $\frac49$, so
the tail after $N$ terms is at most
\[
  2\tau_N + \frac95\Bigl(\frac4K\Bigr)^{n_0}\frac1{n_0^2}, \qquad n_0 = \lceil x\rceil + 1.
\]
About $(p - (K-4)U\log_2e)/\log_2(K/4)$ terms are needed.
\item \emph{Taylor stages.} The truncated series are partial sums of the Taylor series of $Z_t(c_0+\delta)$,
convergent for $|\delta| < c_0$. With $\rho = |\delta|/c_0$, the remainder after $L \ge 2$ terms of every state
component is at most $2e^{4U}c_0^{-1/2}\rho^L/(1-\rho)$. Using
\[
  C_s(c_0) = c_0^{1-s}\,\Gamma\Bigl(s-1, \frac{c_0}U\Bigr), \qquad \Gamma(a, y) \le \frac{y^{a-1}e^{-y}}{1 - (a-1)/y},
\]
one also gets a bound proportional to $E$, of the form
\[
  8e^{4U}E\,\frac{(|\delta|/U)^L}{L!}\Bigl(\frac{U^{1/2}}{c_0} + \frac{U^{3/2}}{c_0^2}\Bigr)
  + \frac{2e^{4U}\rho^{\lfloor c_0/(2U)\rfloor+1}}{c_0^{1/2}(1-\rho)},
\]
so the stages only need to resolve the $p - (c/U)\log_2 e$ significant bits of $J_m$.
\item \emph{Starting values.} With $b_n = 4^n/n!$ we have $|b_n\alpha_n| = (16c)^n/(2n)!$, and
\[
  E\beta_n = D_{n-1/2} - \alpha_nD_{-1/2}, \qquad 0 < D_{n-1/2} \le \frac{EU^{n-1/2}}{n-\frac12}.
\]
All weights are at most $1$, so once
\[
  \frac{16c}{(2N+1)(2N+2)} \le \frac12 \qquad\text{and}\qquad \frac{4U}{N+1} \le \frac12,
\]
every tail of the $\alpha$- and
$\beta$-series is at most twice its first term, which gives an explicit bound of the form
\[
  4D_{-1/2}\frac{(16c)^N}{(2N)!} + O\Bigl(E\,b_NU^{N-1/2}\Bigl(\frac1U + \frac1c\Bigr)\Bigr).
\]
Each $D_{n-1/2}$ splits into parts larger by up to $e^{c/U}$, but in the weighted sums the terms are only
$e^{O(\sqrt c)}$ times the result ($|b_n\alpha_n| \le e^{4\sqrt c}$), so the bound has the size of the result up to such a factor.
\item \emph{The point.} The last stage evaluates at $c_\ell$ with $|c - c_\ell| \le \delta = 2^{-B} + \operatorname{rad}(c)$, and
\[
  |J'(\xi)| \le \frac{2\sqrt\pi\,\Eip(4U)\,e^{-\xi/U}}{U^{1/2}\,\xi} \qquad (\xi \ge c - \delta),
\]
so $B \approx p - (c/U)\log_2e$ bits of $c$ suffice. (Using $\xi \ge c - 1$ instead would lose a factor
$e^{1/U}$: $370$ bits at $U = 4^{-4}$.)
\item \emph{Main term and constants.} The terms of $M$ are at most $3(8U)^N/N!$, and the tail of the series for
$\Li_2(\frac19)$ after $N$ terms is at most $9^{-N}$.
\end{itemize}

\subsection{The logarithmic form}\label{sec:logform}
Theorem~\ref{thm:main} applies directly to
\[
  f(x) = -\log(1-x)\log(1+x) = \sum_{n=1}^{\infty}\frac{A_n}{n}x^{2n}, \qquad A_n = \sum_{j=1}^{2n-1}\frac{(-1)^{j+1}}j,
\]
with $\chi = 1$ and $k_0 = 2$. The kernel is $H(u) = \sum_{n=0}^{\infty}\frac{A_{n+1}}{(n+1)!}u^n$, entire of exponential
type $1$, with
\[
  uH(u) = \log 2\,(e^u-1) + \int_0^1\frac{e^{\beta u}-1}{2\beta(1+\sqrt\beta)}\dd\beta
\]
(from $A_n = \log 2 + \int_0^1 t^{2n-1}/(1+t)\dd t$). Its coefficients satisfy
\[
  (n+1)^2(n+2)(2n+3)h_{n+1} = (n+1)(4n^2+6n+3)h_n - n(2n+1)h_{n-1},
\]
and $H$ satisfies the order-$4$ equation
\[
  \vartheta^2(\vartheta+1)(2\vartheta+1)H - u(\vartheta+1)(4\vartheta^2+6\vartheta+3)H
  + u^2(\vartheta+1)(2\vartheta+3)H = -\frac u2.
\]
The splitting uses $\theta_3$; the $k$-terms are
\[
  I_k = e^{-k^2u_0}\sum_{n=0}^{\infty}\frac{A_{n+1}}{n+1}\frac{e_{n+1}(k^2u_0)}{k^{2n+2}}
\]
(a $5\times5$ triangular recurrence), and the moments $Z_t[H](c)$, from which $J_m = \sqrt\pi Z_{1/2}[H](\pi^2m^2)$,
satisfy a scalar
recurrence of order $6$ in $t$ with coefficients of degree up to $4$ in $c$ and five forcing constants
$(\vartheta^iH)(u_0)$, $i < 4$, and $1$. (As a first-order system in $c$ as in Lemma~\ref{lem:moments} it is linear in $c$,
but of dimension about $16$ instead of $4$.)

This is less efficient than the dilogarithm form for a structural reason: the alternating harmonic numbers
$A_n$ are not hypergeometric (they are a first derivative with respect to a parameter, $\psi(2n) - \psi(n)$),
so the kernel has order $4$ instead of $1$, and the dual system has order $6$ plus forcing instead of $4$, with
quartic instead of linear coefficients; the leaves of the Taylor stages are correspondingly larger. The
dilogarithm form pays for its hypergeometric coefficients with exponential type $4$ instead of $1$ (slower
$k$-series, $k = 2$ excluded), which costs far less. Measured with comparable implementations (CPU time, totals at the
optimal $u_0$):
\begin{center}
\begin{tabular}{lrrrrrr}
\toprule
 & \multicolumn{2}{c}{one $k$-term} & \multicolumn{2}{c}{one dual term} & \multicolumn{2}{c}{total} \\
digits & log & dilog & log & dilog & log & dilog \\
\midrule
$10^4$ & $0.018$\,s & $0.0073$\,s & $0.71$\,s & $0.19$\,s & $19$\,s & $4.9$\,s \\
$10^5$ & $0.25$\,s & $0.078$\,s & $18.5$\,s & $3.4$\,s & $1\,165$\,s & $280$\,s \\
$10^6$ & $\approx 4$\,s & $1.0$\,s & $300$\,s & $56$\,s & $\approx 59\,000$\,s & $12\,880$\,s \\
\bottomrule
\end{tabular}
\end{center}
(both before the optimizations of Section~\ref{sec:khimpl}; the dilogarithm form is $3.9$--$4.6$ times faster, with
the ratio slowly growing with the precision). The second identity of \cite[\S3]{BBC97},
\[
  \log K\log2 = \frac{\pi^2}{6} - \frac12\log^22 + \sum_{k=2}^{\infty}\Li_2\Bigl(-\frac1{k^2-1}\Bigr)
\]
(from Abel's dilogarithm identity), has a
kernel $-e^u\Ein(u)/u$ of type $1$, but a five-dimensional dual state; we did not implement it.

\subsection{Implementation}\label{sec:khimpl}
Our implementation is part of FLINT \cite{FLINT}, as the function
\texttt{mp\_real\_const\_khinchin} of the \texttt{mp\_real} module. It uses the theta splitting above
$5\cdot10^6$ bits (about $1.5\cdot10^6$ digits) and an optimized zeta-value expansion \eqref{eq:zetaseries} below.
We only sketch it.
\begin{itemize}
\item All linear recurrences (the $k$-terms, the Taylor stages and starting values of the dual terms,
$\Gamma(\frac12, y)$) are evaluated by exact binary splitting of products of small integer matrices over chunks of
terms whose products reach about half the working precision. Within a chunk, the two halves are merged by one
augmented product with FLINT's FFT-based matrix multiplication (\texttt{fmpz\_mat\_mul\_fft\_small}), which
transforms each entry once. Each chunk product is applied to the state, kept as integers scaled by the denominators,
by one truncated transformed matrix--vector product with per-entry exponents, so that all sizes stay $O(p)$.
\item The $k$-terms use the triangular $2\times2$ recurrence above; the lowest levels of their trees are done with
fixed-size limb arrays.
\item $\Gamma(\frac12, y)$ at the dyadic points $y = c_0/u_0$ is computed by the upper (asymptotic) series when it is
accurate enough, otherwise by binary splitting of the Legendre continued fraction (whose consecutive convergents
enclose the value) or of the lower series, whichever is shorter.
\item The bit-burst uses growth factor $4$ per stage (the stage cost is about $(g+1)/(g-1)$ times a plain stage,
because the coefficients are linear in the previous point); short final stages are summed directly.
\item Each term is computed with only the relative precision it needs (Remark~\ref{rem:practical}), and every omitted
term and Taylor remainder is bounded rigorously (Section~\ref{sec:khbounds}).
\item The $k$-terms are taken in chunks of consecutive $k$, which share powers of $e^{-u_0}$; the dual terms are taken
largest first, with at most half of the threads on dual terms while $k$-chunks remain. The finished terms are added
exactly in fixed point, so that no term is kept and the result does not depend on the scheduling. The memory is linear
in the precision times the number of threads, mostly for FFT transforms.
\end{itemize}

\subsection{Results}\label{sec:khresults}
Table~\ref{tab:khinchin} compares the theta splitting (with $u_0 = \frac1{16}$) with the zeta-value expansion in FLINT,
both run with $8$ threads on the same machine, with an AMD Ryzen 7 PRO 5850U processor and $32$ GB of memory. The
comparison is with a strong baseline: the zeta-value implementation in FLINT has been optimized meticulously. On the
same machine it is $3$ times faster than the program of Ghazi \cite{Gha26}, used for the previous record, at $10^5$
digits and $5$ times faster at $3\cdot10^5$ digits, although that program is built on FLINT as well and uses the same
reverse Bernoulli iterator.
\begin{table}[ht]
\centering\footnotesize
\setlength{\tabcolsep}{3.5pt}
\begin{tabular}{rrrrrrrr}
\toprule
& & \multicolumn{3}{c}{theta splitting} & \multicolumn{3}{c}{zeta values} \\
\cmidrule(lr){3-5}\cmidrule(lr){6-8}
digits & $k$/dual terms & CPU & wall & memory & CPU & wall & memory \\
\midrule
$10^5$ & $1\,919$/$38$ & $88.3$\,s & $11.2$\,s & $223$\,MB & $18.7$\,s & $2.43$\,s & $78$\,MB \\
$10^6$ & $6\,069$/$121$ & $5\,416$\,s & $679$\,s & $1.08$\,GB & $4\,239$\,s & $544$\,s & $1.27$\,GB \\
$2\cdot10^6$ & $8\,583$/$171$ & $18\,815$\,s & $2\,355$\,s & $1.78$\,GB & $23\,485$\,s & $3\,372$\,s & $1.96$\,GB \\
$10^7$ & $19\,194$/$382$ & $327\,183$\,s & $40\,921$\,s & $10.5$\,GB & $1\,852\,090$\,s & $235\,475$\,s & $15.3$\,GB \\
\bottomrule
\end{tabular}
\vspace{\abovecaptionskip}
\caption{Khinchin's constant with $8$ threads: CPU time (summed over the threads), wall time and peak memory.}
\label{tab:khinchin}
\end{table}
Per decade the theta splitting grows by a factor $61$ from $10^5$ to $10^6$ digits and $60$ from $10^6$ to $10^7$ (about
$p^{1.78}$, consistent with $p^{3/2}$ times logarithmic factors), the zeta-value expansion by $227$ and $437$ (about
$p^{2.36}$ and $p^{2.64}$; the exponent of this power-sum based implementation keeps growing with $p$). The crossover
lies between $10^6$ and $2\cdot10^6$ digits, as reflected in the cutoff of \texttt{mp\_real\_const\_khinchin}. At $10^7$
digits the theta splitting is $5.7$ times faster ($11.4$ hours of wall time against $65.4$) and needs less memory
($10.5$ GB against $15.3$ GB). Both methods parallelize almost perfectly: the ratio of CPU to wall time is $7.9$--$8.0$ on
$8$ threads. The $10^7$-digit result of the theta splitting has error radius $1.8\cdot10^{-10\,000\,044}$, and the
two results agree. Since the theta splitting shares nothing with the zeta-value method (no Bernoulli numbers or power
sums), it is a genuinely independent check.

\subsection{The continued fraction of $K$}\label{sec:khcf}
With $D = 10^7$ and $d = \lfloor 10^DK\rfloor$, we have $K \in I = [d\,10^{-D}, (d+1)10^{-D}]$. We computed the
continued fraction expansions of both endpoints of $I$ with FLINT's subquadratic \texttt{fmpq\_get\_cfrac} (in
$16$ seconds). They agree in the first $N + 1$ partial quotients, $N = 9\,699\,141$,
\[
  K = [2; 1, 2, 5, 1, 1, 2, 1, 1, 3, 10, 2, 1, 3, 2, 24, 1, 3, 2, 3, \ldots, a_N, \ldots],
\]
(OEIS A002211 \cite{OEIS}) and differ in $a_{N+1}$, which is $9$ for one endpoint and $26$ for the other. Since the set of
real numbers whose expansion begins with $a_0, \ldots, a_N$ is an interval, these are the first $N + 1$ partial quotients of
$K$. Let $p_n/q_n$ be the convergents. Then $q_N = 2.484246\ldots\cdot10^{4\,999\,999}$, and the largest of these
partial quotients is $a_{7\,384\,031} = 34\,770\,305$.

The statistics of the partial quotients are those of a typical real number. Lochs' theorem \cite{Loc64} says that
for almost all $x$, $n$ decimal digits determine about $\lambda n$ partial quotients, and L\'evy's theorem
\cite{Lev36} that $\frac1n\log q_n \to \mu$, where
\[
  \lambda = \frac{6\log2\log10}{\pi^2} = 0.970270\ldots, \qquad \mu = \frac{\pi^2}{12\log2} = 1.186569\ldots;
\]
here $N/D = 0.969914$ and $\frac1N\log q_N = 1.187005$. The proportions of the partial quotients $1$, $2$ and $3$ are
$0.414921$, $0.169996$ and $0.093136$, against the Gauss--Kuzmin probabilities $\log_2\frac43 = 0.415037$, $0.169925$
and $0.093109$.

\begin{theorem}\label{thm:khrat}
If $K = p/q$ with integers $p$ and $q > 0$, then $q \ge 10q_N + q_{N-1} > 2.587\cdot10^{5\,000\,000}$, and
$p > 6.94\cdot10^{5\,000\,000}$.
\end{theorem}

\begin{proof}
Let
\[
  x_t = \frac{tp_N + p_{N-1}}{tq_N + q_{N-1}} = [a_0; a_1, \ldots, a_N, t].
\]
Its denominator is $q(x_t) = tq_N + q_{N-1}$ (in lowest terms), and consecutive $x_t$ are Farey neighbours, since
\begin{multline*}
  \bigl((t+1)p_N + p_{N-1}\bigr)(tq_N + q_{N-1}) - (tp_N + p_{N-1})\bigl((t+1)q_N + q_{N-1}\bigr) \\
  = p_Nq_{N-1} - p_{N-1}q_N = \pm1.
\end{multline*}
Exact rational arithmetic shows that $I$ lies
between $x_9$ and $x_{27}$ and contains $x_{10}, \ldots, x_{26}$. A rational number strictly between two Farey
neighbours $a/b$ and $c/e$ has denominator at least $b + e$. Hence every rational number in $I$ is either one of
$x_{10}, \ldots, x_{26}$ or has denominator at least $q(x_9) + q(x_{10})$; in both cases its denominator is at least
$q(x_{10}) = 10q_N + q_{N-1}$. The bound for $p$ follows from $K > 2.685$.
\end{proof}

The theorem depends, of course, on the correctness of the computed digits, which rests on the rigorous error bounds
of Section~\ref{sec:khbounds} and on the agreement with the zeta-value computation.

Khinchin's theorem \cite{Khi35} says that the geometric mean
\[
  G_n = (a_1a_2\cdots a_n)^{1/n}
\]
of the partial quotients tends to $K$ for almost all $x$. Whether $K$ itself has this property is not known, but its first $N$ partial
quotients behave as expected (Figure~\ref{fig:khcf}): $G_N = 2.686720\ldots$, so that $G_N - K = 1.27\cdot10^{-3}$. The
error decreases like $n^{-1/2}$, as it should: for $x$ chosen at random (with respect to the Gauss or the Lebesgue
measure) the sums $\sum_{k=1}^{n}\log a_k$ satisfy a central limit theorem \cite{IK02}, so that
\[
  z_n = \sqrt n\,(\log G_n - \log K)
\]
is asymptotically normal with mean $0$ and variance
\[
  \sigma^2 = \lim_{n\to\infty}\frac1n\operatorname{Var}\sum_{k=1}^{n}\log a_k;
\]
this is smaller than
$\operatorname{Var}(\log a_1) = 1.3919$, because consecutive partial quotients are weakly correlated. Batch means over
blocks of $10^4$ consecutive partial quotients of $K$ give $\sigma^2 \approx 1.21 \pm 0.06$, that is,
$\sigma \approx 1.10$. In the time scale $\log n$ the normalized error $z_n$ wanders like a Brownian motion, as the
invariance principle predicts; it stays below $3.6$ in absolute value for $n \ge 10^3$ (with excursions to $-3.57$ at
$n = 1\,289$ and $3.21$ at $n = 4\,386\,910$, unremarkable over so long a range of $\log n$), and
$z_N = 1.47 \approx 1.34\,\sigma$.

\begin{figure}[ht]
\centering
\begin{tikzpicture}
\begin{axis}[
  name=gm, width=0.46\textwidth, height=0.42\textwidth,
  xmode=log, xmin=10, xmax=1.2e7, ymin=1.9, ymax=3.5,
  xlabel={$n$}, ylabel={$G_n$},
  ytick={2.0,2.5,3.0,3.5}, yticklabel style={/pgf/number format/fixed, /pgf/number format/precision=1},
  tick label style={font=\footnotesize}, label style={font=\small},
  legend style={font=\footnotesize, at={(0.97,0.97)}, anchor=north east, draw=none, fill=none}]
\addplot[gray!60, thin, domain=10:1.2e7, samples=200, forget plot] {2.685452001*exp(2*1.10/sqrt(x))};
\addplot[gray!60, thin, domain=10:1.2e7, samples=200, forget plot] {2.685452001*exp(-2*1.10/sqrt(x))};
\addplot[black, dashed, domain=10:1.2e7, samples=2] {2.685452001};
\addlegendentry{$K$}
\addplot[blue!70!black, thick] table[x=n, y=gm] {khcf_gm.dat};
\addlegendentry{$G_n$}
\end{axis}
\begin{axis}[
  at={(gm.east)}, anchor=west, xshift=0.8cm, width=0.46\textwidth, yticklabel pos=right, ylabel near ticks, height=0.42\textwidth,
  xmode=log, xmin=10, xmax=1.2e7, ymin=-4.2, ymax=4.2,
  xlabel={$n$}, ylabel={$z_n$},
  ytick={-4,-2,0,2,4},
  tick label style={font=\footnotesize}, label style={font=\small}]
\addplot[gray!60, thin, domain=10:1.2e7, samples=2] {2*1.10};
\addplot[gray!60, thin, domain=10:1.2e7, samples=2] {-2*1.10};
\addplot[black, dashed, domain=10:1.2e7, samples=2] {0};
\addplot[blue!70!black] table[x=n, y=z] {khcf_gm.dat};
\end{axis}
\end{tikzpicture}
\caption{The geometric mean $G_n$ of the first $n$ partial quotients of $K$ (left) and the normalized error
$z_n = \sqrt n\,(\log G_n - \log K)$ (right), for $10 \le n \le N = 9\,699\,141$, at about $1\,100$ logarithmically spaced $n$. The grey curves show
$K e^{\pm2\sigma/\sqrt n}$ (left) and $\pm2\sigma$ (right), with $\sigma = 1.10$ estimated from the partial quotients.}
\label{fig:khcf}
\end{figure}

%======================================================================
\section{Other constants}\label{sec:apps}

We applied Algorithm~\ref{alg:main} to the following constants with a generic program, \texttt{holsum}
(\url{https://github.com/fredrik-johansson/holonomic-sums}), which is much less optimized than the implementation of Section~\ref{sec:khimpl}: ball arithmetic
throughout, FLINT's default matrix products, no fixed-point levels, heuristic (but conservative) truncation of the
series, the same growth factor $4$, and split point $u_0 = \frac1{16}$ for the constants of
Sections~\ref{sec:AG} and~\ref{sec:KH} and $u_0 = \frac14$ for the other two. (The program adds a margin of two terms
to the counts \eqref{eq:counts}.) The constants of Sections~\ref{sec:AG} and~\ref{sec:KH} have $\chi = 1$ and $k_0 = 2$, so
\[
  \Theta(u) = \sum_{k=2}^{\infty}e^{-k^2u} = \frac12\sqrt{\frac\pi u} - \frac12 - e^{-u} + \sqrt{\frac\pi u}\sum_{m=1}^{\infty}e^{-\pi^2m^2/u}
\]
and
\[
  S = \sum_{k=2}^{\infty} I_k + \sqrt\pi\sum_{m=1}^{\infty}J_m + M, \qquad
  M = \int_0^{u_0}\bigl(\tfrac12\sqrt{\pi/u} - \tfrac12 - e^{-u}\bigr)h(u)\dd u,
\]
with $J_m = Z_{1/2}[h_0](\pi^2m^2) + Z_1[h_1](\pi^2m^2)$.

\subsection{The Alladi--Grinstead and L\"uroth constants}\label{sec:AG}
The constant
\[
  \mathcal A = \sum_{k=2}^{\infty}\frac{\log k}{k^2-1} = 1.02313872642793929553\ldots
\]
(OEIS A340440; twice it is A244109, related to L\"uroth representations \cite{Kal88,Fin03}) equals
$-\sum_{i=1}^{\infty}\zeta'(2i)$ and also
\[
  \mathcal A = \tfrac12\log2 + \tfrac12 S_A, \qquad
  S_A = \sum_{k=2}^{\infty}\frac{2}{k}\arctanh\frac1k = \sum_{m=1}^{\infty}\frac{2(\zeta(2m)-1)}{2m-1}
\]
(the second form by summation by parts, or from the zeta-value expansion). For $S_A$, we have
$f(x) = 2x\arctanh x$ and $a_{2m} = \frac{2}{2m-1}$, so that
\[
  h_0(u) = \sum_{j=0}^{\infty}\frac{2u^j}{(2j+1)\,j!} = \sqrt{\pi/u}\,\erfi\sqrt u, \qquad \vartheta h_0 = e^u - \tfrac12h_0 .
\]
The dual moments $\nu_t = Z_t[h_0]$, $Y_t = Z_t[e^u]$ satisfy
\begin{align*}
  c\,\nu_{t+1} &= Eu_0^{1-t}h_0(u_0) - (\tfrac12 - t)\nu_t - Y_t, \\
  c\,Y_{t+1} &= Eu_0^{1-t}e^{u_0} - (1-t)Y_t - Y_{t-1},
\end{align*}
a state of dimension $3$ ($\nu_t, Y_t, Y_{t-1}$) with two forcing constants, $t_0 = \frac12$. The $k$-terms are
\[
  I_k = e^{-x}\sum_{m=1}^{\infty} \frac{2}{2m-1}\,\frac{e_m(x)}{K^m},
\]
a $2\times2$ recurrence as for Khinchin's constant.

Next, let
\begin{align*}
  \mathcal C &= \sum_{k=2}^{\infty}\frac1k\log\frac{k}{k-1} = \sum_{k=1}^{\infty}\frac{\log k}{k(k+1)} = \sum_{n=1}^{\infty}\frac{\zeta(n+1)-1}{n} \\
  &= 0.78853056591150896106\ldots
\end{align*}
(A085361). Then
\begin{align*}
  e^{\mathcal C-1} &= 0.80939402054063913071\ldots, \\
  e^{\mathcal C} &= \prod_{k=1}^{\infty}k^{1/(k(k+1))} = 2.20016105809902655319\ldots
\end{align*}
are the Alladi--Grinstead constant \cite{AG77} (A085291) and Kalpazidou's analogue of Khinchin's constant for
L\"uroth expansions \cite{Kal88} (A245254). With
$f(x) = -x\log(1-x) = x\arctanh x - \frac x2\log(1-x^2)$ we have $\mathcal C = \frac12S_A + \mathcal O$, where
\[
  \mathcal O = -\sum_{k=2}^{\infty}\frac{1}{2k}\log\Bigl(1-\frac1{k^2}\Bigr) = \frac12\sum_{m=1}^{\infty}\frac{\zeta(2m+1)-1}{m},
\]
with $a_{2m+1} = \frac1{2m}$. The odd kernel and an auxiliary function are
\[
  h_1(u) = \sum_{m=1}^{\infty}\frac{u^m}{2m\,\Gamma(m+\frac12)}, \qquad
  \varphi(u) = \sum_{m=1}^{\infty}\frac{u^m}{\Gamma(m+\frac12)} = \sqrt u\,e^u\erf\sqrt u,
\]
with $\vartheta h_1 = \frac12\varphi$ and $\vartheta\varphi = (u+\frac12)\varphi + u/\sqrt\pi$. With $\nu_t = Z_t[h_1]$ and
$\Phi_t = Z_t[\varphi]$,
\begin{align*}
  c\,\nu_{t+1} &= Eu_0^{1-t}h_1(u_0) - (1-t)\nu_t - \tfrac12\Phi_t, \\
  c\,\Phi_{t+1} &= Eu_0^{1-t}\varphi(u_0) - (\tfrac32-t)\Phi_t - \Phi_{t-1} - C_{t-1}/\sqrt\pi,
\end{align*}
a state of dimension $4$ ($\nu_t, \Phi_t, \Phi_{t-1}, C_{t-1}$) with three forcing constants and $t_0 = 1$, so that the
starting values involve $E_1(c_0/u_0)$. The $k$-terms are, by \eqref{eq:Ik},
\[
  \frac1k\biggl[-\frac12\log\Bigl(1-\frac1{k^2}\Bigr)\erfc(k\sqrt{u_0})
  + \frac{2}{\sqrt\pi}e^{-x}\sqrt x\sum_{m=1}^{\infty} \frac{K^{-m}}{2m}\sum_{j=0}^{m-1}\frac{(2x)^j}{3\cdots(2j+1)}\biggr].
\]
One run computes $S_A$ and $\mathcal O$ together.

\subsection{Khinchin's harmonic mean}\label{sec:KH}
For almost all $x$, Khinchin's H\"older means of the partial quotients satisfy, for $p < 1$ \cite{BBC97},
\[
  K_p = \lim_{n\to\infty}\Bigl(\frac1n\sum_{i=1}^{n} a_i^p\Bigr)^{1/p}, \qquad
  K_p^p = \sum_{k=1}^{\infty}k^p\log_2\Bigl(1 + \frac1{k(k+2)}\Bigr).
\]
For $p = -1$, with $j = k+1$,
\[
  K_{-1} = \frac{\log 2}{T}, \qquad T = \sum_{j=2}^{\infty}\frac{-\log(1-1/j^2)}{j-1},
\]
so that $f(x) = -\frac{x}{1-x}\log(1-x^2)$ and $a_n = H_{\lfloor(n-1)/2\rfloor}$.
Both kernels involve harmonic numbers:
\begin{align*}
  h_0(u) &= \sum_{j=0}^{\infty}\frac{H_j u^j}{j!} = e^u\Ein(u), & \vartheta h_0 &= uh_0 + e^u - 1, \\
  h_1(u) &= \sum_{m=1}^{\infty}\frac{H_mu^m}{\Gamma(m+\frac12)}, & \vartheta h_1 &= (u+\tfrac12)h_1 - \eta + \varphi,
\end{align*}
with $\eta$, $\varphi$ the functions $h_1$ and $\varphi$ of Section~\ref{sec:AG}. With $Y_t = Z_t[e^u]$, the even
and odd dual states are
\[
  (\nu_t, \nu_{t-1}, Y_t, Y_{t-1}, C_t), \qquad (\nu_t, \nu_{t-1}, \eta_t, \Phi_t, \Phi_{t-1}, C_{t-1}),
\]
and
the $k$-side recurrences have states of dimension $4$ (the harmonic numbers double the $2\times2$ systems). The
constant $\sum_{m=1}^{\infty} H_mK^{-m} = -\log(1-\frac1K)/(1-\frac1K)$ multiplies $\erfc(k\sqrt{u_0})$.
The same applies to every $K_{-n}$, $n \ge 1$, with $f(x) = (x/(1-x))^n(-\log(1-x^2))$.

\subsection{The Hardy--Littlewood function}\label{sec:HL}
The function
\[
  H(x) = \sum_{k=1}^{\infty}\frac1k\sin\frac xk,
\]
studied by Hardy and Littlewood, is a standard test case for the summation of slowly convergent series
\cite{Gau05,Mon10}. Here $f(t) = t\sin(xt)$ is entire, so $k_0 = 1$, and the
summand $f(1/k)$ has an essential singularity at $k = 0$. For integer $x$ the kernel is
\[
  h(u) = \sum_{j=0}^{\infty}\frac{(-1)^jx^{2j+1}u^j}{(2j+1)!\,j!} = x\,{}_0F_2\Bigl(;\frac32,1;-\frac{x^2u}4\Bigr),
\]
with $\vartheta^2(\vartheta+\frac12)h = -\frac{x^2}{4}u\,h$,
an entire function of order $\frac13$, so the $k$-side series converge for every $k \ge 1$. With
$P_t = Z_t[h]$, $Q_t = Z_t[\vartheta h]$ and $R_t = Z_t[\vartheta^2h]$, \eqref{eq:ibp} gives
\begin{align*}
  c\,P_{t+1} &= E\,u_0^{1-t}\,h(u_0) - (1-t)P_t - Q_t, \\
  c\,Q_{t+1} &= E\,u_0^{1-t}\,(\vartheta h)(u_0) - (1-t)Q_t - R_t, \\
  c\,R_{t+1} &= E\,u_0^{1-t}\,(\vartheta^2h)(u_0) - \bigl(\tfrac12 - t\bigr)R_t + \tfrac{x^2}{4}P_{t-1},
\end{align*}
and $J_m = Z_{1/2}[h](\pi^2m^2)$: a state of dimension $4$ ($P_t, Q_t, R_t, P_{t-1}$) with three forcing constants and
$t_0 = \frac12$, so that the starting values involve $\Gamma(\frac12, c_0/u_0)$. The $k$-terms are
\[
  I_k = e^{-x_k}\sum_{m=1}^{\infty}\frac{(-1)^{m-1}x^{2m-1}}{(2m-1)!}K^{-m}e_m(x_k),
\]
a $2\times2$ recurrence as for Khinchin's constant.

\subsection{Periodic arrays of vortex rings}\label{sec:vortex}

Consider an infinite row of identical coaxial thin-cored vortex rings of radius $r$ and circulation $\Gamma$, spaced
$\lambda r$ apart along their common axis. An isolated thin ring travels along its axis with a self-induced velocity
that depends logarithmically on the radius of its core (Kelvin's formula). In the row, every ring is moreover carried
along by the flow that all the other rings induce at its core. By the Biot--Savart law, a ring at axial distance
$n\lambda r$ induces at a point of the core of another ring the axial velocity
$\frac{\Gamma}{4\pi r}\int_0^{2\pi}(1-\cos\varphi)\,(2(1-\cos\varphi)+n^2\lambda^2)^{-3/2}\dd\varphi$, where $\varphi$ is the
azimuthal angle along the inducing ring, measured from the point. This is positive: the ring is pushed forward
alike by the rings ahead of it and by those behind it, whose radial contributions cancel, so the whole row translates
rigidly. Summing over the rings on both sides, the speed of the row exceeds that of an isolated ring by
$U(\lambda)\,\Gamma/r$, where
\begin{equation}\label{eq:vortex}
  U(\lambda) = \frac{1}{2\pi}\sum_{n=1}^{\infty}\int_0^{2\pi}\frac{(1-\cos\varphi)\dd\varphi}{\bigl(2(1-\cos\varphi)+n^2\lambda^2\bigr)^{3/2}},
\end{equation}
a formula found independently by Vasilev (1916) and by Levy and Forsdyke (1928), see \cite{MS22}. It is a sum of the
form \eqref{eq:S}: the summand is a complete elliptic integral, and \cite{MS22} sums it by truncating after $500$ terms.

\subsubsection{The series}
Substituting $s = \sin^2(\varphi/2)$, expanding the integrand binomially and using Wallis's integrals gives, for $a > 2$,
\[
  \frac1{2\pi}\int_0^{2\pi}\frac{(1-\cos\varphi)\dd\varphi}{(2(1-\cos\varphi)+a^2)^{3/2}}
  = a^{-3}\,{}_2F_1\Bigl(\frac32,\frac32;2;-\frac4{a^2}\Bigr),
\]
and the identity extends to all $a > 0$ by analytic continuation. By contiguity, ${}_2F_1(\frac32,\frac32;2;z)$ is a
combination of the complete elliptic integrals $K$ and $E$ with coefficients rational in $z$. Hence, with
$\beta = 4/\lambda^2$,
\[
  U(\lambda) = \lambda^{-3}\sum_{n=1}^{\infty} n^{-3}\,F\Bigl(-\frac{\beta}{n^2}\Bigr), \qquad F(z) = {}_2F_1\Bigl(\frac32,\frac32;2;z\Bigr).
\]
This is $\sum_{n=1}^{\infty} f(1/n)$ with $f(x) = \lambda^{-3}x^3F(-\beta x^2)$, a holonomic function with radius of convergence
$\rho = \lambda/2$. Theorem~\ref{thm:main} applies with $k_0 > 2/\lambda$, the terms $n \le 2/\lambda$ being added
separately, so $U(\lambda)$ can be computed to $p$ bits in time $\widetilde O(p^{3/2})$ for every rational $\lambda$ (and,
with algebraic $\beta$, for every algebraic $\lambda$). The leading term gives $U(\lambda) \sim \zeta(3)\lambda^{-3}$ as
$\lambda \to \infty$ \cite{MS22}. The zeta-value expansion of the terms $n \ge k_0$,
\[
  \lambda^{-3}\sum_{j=0}^{\infty}\frac{(3/2)_j^2}{(2)_j\,j!}(-\beta)^j\Bigl(\zeta(2j+3) - \sum_{n=1}^{k_0-1}n^{-2j-3}\Bigr),
\]
converges like $(\beta/k_0^2)^j$.

\subsubsection{The splitting}
The kernel has a closed form. Let
\[
  y(u) = {}_1F_1\Bigl(\frac32; 2; -\beta u\Bigr) = e^{-\beta u/2}\Bigl(I_0\bigl(\tfrac{\beta u}{2}\bigr) - I_1\bigl(\tfrac{\beta u}2\bigr)\Bigr).
\]
Termwise integration gives
\[
  n^{-3}F\Bigl(-\frac\beta{n^2}\Bigr) = \frac{2}{\sqrt\pi}\int_0^\infty u^{1/2}y(u)e^{-n^2u}\dd u.
\]
Since
$y(u) \sim (\beta u)^{-3/2}/\sqrt\pi$ as $u \to +\infty$, the kernel $u^{1/2}y(u)$ is bounded on $(0,\infty)$, and this
representation holds for every $n \ge 1$, not only for $n > 2/\lambda$. Splitting at $u_0$ as in
Proposition~\ref{prop:split}, with
\[
  \sum_{n=1}^{\infty}e^{-n^2u} = \frac12\Bigl(\sqrt{\frac\pi u} - 1\Bigr) + \sqrt{\frac\pi u}\sum_{m=1}^{\infty}e^{-\pi^2m^2/u},
\]
gives
\[
  U(\lambda) = \frac{2}{\sqrt\pi\,\lambda^3}\Bigl(\sum_{n=k_0}^{\infty}I_n + \sum_{n=1}^{k_0-1}(T_n - L_n) + M + \sqrt\pi\sum_{m=1}^{\infty}J_m\Bigr),
\]
where, with $K = n^2$ and $x = Ku_0$,
\begin{align*}
  I_n &= \int_{u_0}^\infty u^{1/2}y\,e^{-Ku}\dd u
       = \frac{1}{2K^{3/2}}\Bigl(\Gamma(\tfrac12, x)\,F\bigl(-\tfrac\beta K\bigr)
         + 2e^{-x}\sqrt x\sum_{j=0}^{\infty}A_jR_j\Bigr), \\
  &\qquad A_j = \frac{(3/2)_j^2}{(2)_j\,j!}\Bigl(-\frac\beta K\Bigr)^j, \quad R_j = \sum_{i=0}^{j}\frac{\Gamma(\frac32)\,x^i}{\Gamma(i+\frac32)}, \\
  T_n &= \int_0^\infty u^{1/2}y\,e^{-Ku}\dd u = \frac{\sqrt\pi}{2}(K+\beta)^{-3/2}\,{}_2F_1\Bigl(\frac32,\frac12;2;\frac{\beta}{\beta+K}\Bigr), \\
  L_n &= \int_0^{u_0}u^{1/2}y\,e^{-Ku}\dd u, \\
  M &= \frac{\sqrt\pi}2\int_0^{u_0}y\dd u - \frac12\int_0^{u_0}u^{1/2}y\dd u, \\
  J_m &= \int_0^{u_0}y(u)\,e^{-\pi^2m^2/u}\dd u .
\end{align*}
The $k$-side series converge like $(\beta/K)^j$, so we take $k_0$ minimal with $k_0^2 > 2\beta$; the finitely many terms
$n < k_0$ are the full integral $T_n$ (by Pfaff's transformation) minus $L_n$, whose integrand $u^{1/2}y(u)e^{-Ku}$ has an
entire power series factor with coefficients satisfying a three-term recurrence. $F(-\beta/K) = \sum_{j=0}^{\infty}A_j$ is summed in the
same binary splitting as $\sum_{j=0}^{\infty}A_jR_j$, with the state $(A_jR_j, A_jx^j/\Gamma(j+\frac32), A_j)$.

\subsubsection{Dual side} With $\vartheta y = w$, the equation of $y$ is $\vartheta w = -w - \beta u(w + \frac32 y)$. For
$P_t = Z_t[y]$ and $Q_t = Z_t[w]$ (Section~\ref{sec:dualgen}), \eqref{eq:ibp} gives
\begin{align*}
  c\,P_{t+1} &= E\,u_0^{1-t}\,y(u_0) - (1-t)P_t - Q_t, \\
  c\,Q_{t+1} &= E\,u_0^{1-t}\,w(u_0) + t\,Q_t + \beta Q_{t-1} + \tfrac{3\beta}{2}P_{t-1},
\end{align*}
and $J_m = P_0(\pi^2m^2)$: a system of order $4$ (state $P_t, Q_t, P_{t-1}, Q_{t-1}$) with coefficients linear in $c$ and
two forcing constants. The moment indices are integers, so the starting values at $c_0$ are combinations of
$E_1(c_0/u_0)$ and $e^{-c_0/u_0}$: with $y = \sum_{j=0}^{\infty}y_ju^j$ and $D_s = \alpha_sE_1(c_0/u_0) + \beta_sE$ as in
Section~\ref{sec:dualgen},
\begin{align*}
  P_0 &= \sum_{j=0}^{\infty}y_jD_{j+1}, & Q_0 &= \sum_{j=0}^{\infty} jy_jD_{j+1}, \\
  P_{-1} &= \sum_{j=0}^{\infty}y_jD_{j+2}, & Q_{-1} &= \sum_{j=0}^{\infty}jy_jD_{j+2},
\end{align*}
all summed by one binary splitting with a three-term state.

For small $\lambda$ the sum samples the function $a \mapsto a^{-3}F(-4/a^2)$ at the fine grid $a = n\lambda$, and the method
becomes more expensive: there are about $2.8/\lambda$ terms $n < k_0$, and the power series of $y$ at $u_0$, which
enter the main term and the starting values, have terms of size up to $e^{\beta u_0}$, so that about $5.8\,u_0/\lambda^2$ guard
bits are needed. For $\lambda \ge p^{-1/2}$ these costs stay within the bound of Theorem~\ref{thm:main}.

\subsection{Computations}\label{sec:appcomp}
Tables~\ref{tab:apps} and~\ref{tab:appsref} show timings of \texttt{holsum} with $8$ threads (on the machine of
Section~\ref{sec:khresults}) for $\mathcal A$, $\mathcal C$, $T$ (for $K_{-1}$), the vortex sum at $\lambda = 1$
($\beta = 4$, $k_0 = 3$) and $H(1)$. Besides the theta splitting, the program implements two reference methods. The
zeta-value expansion \eqref{eq:zetaseries} is applied after taking out the terms $k \le 31$, so that it converges
like $32^{-n}$, except for
\[
  H(x) = \sum_{m=0}^{\infty}(-1)^m\zeta(2m+2)\frac{x^{2m+1}}{(2m+1)!},
\]
whose coefficients decay factorially; the zeta values come from FLINT, in blocks of consecutive arguments. Euler--Maclaurin summation uses the
cutoff $M \approx p/4$ for $p$ bits; the Taylor coefficients of the summand at $M$ come from a linear recurrence (from
its differential equation), the Bernoulli numbers from FLINT's reverse iterator. All results agree to all digits
where they overlap, and the $10^6$-digit values agree with the $10^5$-digit ones. At $10^5$ digits, the theta
splitting also agrees with runs at a second split point, which change every $k$-term, dual term and the main term.

\begin{table}[tp]
\centering\small
\begin{tabular}{llrrrr}
\toprule
constant & digits & $k$/dual terms & CPU ($k$-side) & wall & memory \\
\midrule
$\mathcal A$ & $10^4$ & $606$/$12$ & $2.35$\,s ($0.52$) & $0.39$\,s & $58$\,MB \\
 & $10^5$ & $1\,918$/$38$ & $137$\,s ($0.58$) & $17.9$\,s & $247$\,MB \\
 & $10^6$ & $6\,068$/$120$ & $8\,435$\,s ($0.55$) & $1\,063$\,s & $1.43$\,GB \\
\midrule
$\mathcal C$ & $10^4$ & $606$/$12$ & $7.00$\,s ($0.66$) & $1.01$\,s & $64$\,MB \\
 & $10^5$ & $1\,918$/$38$ & $469$\,s ($0.68$) & $59.7$\,s & $368$\,MB \\
 & $10^6$ & $6\,068$/$120$ & $29\,201$\,s ($0.66$) & $3\,662$\,s & $2.11$\,GB \\
\midrule
$T$ ($K_{-1}$) & $10^4$ & $606$/$12$ & $14.7$\,s ($0.58$) & $2.05$\,s & $88$\,MB \\
 & $10^5$ & $1\,918$/$38$ & $997$\,s ($0.59$) & $127$\,s & $560$\,MB \\
 & $10^6$ & $6\,068$/$120$ & $51\,551$\,s ($0.55$) & $6\,458$\,s & $4.01$\,GB \\
\midrule
$U(1)$ & $10^4$ & $303$/$24$ & $9.85$\,s ($0.52$) & $1.34$\,s & $73$\,MB \\
 & $10^5$ & $959$/$76$ & $677$\,s ($0.56$) & $85.6$\,s & $305$\,MB \\
 & $10^6$ & $3\,034$/$241$ & $37\,795$\,s ($0.50$) & $4\,736$\,s & $2.19$\,GB \\
\midrule
$H(1)$ & $10^4$ & $303$/$24$ & $6.15$\,s ($0.35$) & $0.86$\,s & $73$\,MB \\
 & $10^5$ & $959$/$76$ & $419$\,s ($0.41$) & $52.8$\,s & $419$\,MB \\
 & $10^6$ & $3\,034$/$241$ & $28\,465$\,s ($0.36$) & $3\,564$\,s & $2.44$\,GB \\
\bottomrule
\end{tabular}
\vspace{\abovecaptionskip}
\caption{Timings of the theta splitting in \texttt{holsum} with $8$ threads: CPU time summed over the threads, with
the fraction spent on the $k$-side in parentheses, wall time and peak memory.}
\label{tab:apps}

\bigskip

\begin{tabular}{llrrrr}
\toprule
& & \multicolumn{2}{c}{zeta values} & \multicolumn{2}{c}{Euler--Maclaurin} \\
\cmidrule(lr){3-4}\cmidrule(lr){5-6}
constant & digits & CPU & wall & CPU & wall \\
\midrule
$\mathcal A$ & $10^4$ & $4.26$\,s & $0.81$\,s & $2.51$\,s & $0.82$\,s \\
 & $10^5$ & $1\,059$\,s & $134$\,s & $1\,695$\,s & $353$\,s \\
\midrule
$\mathcal C$ & $10^4$ & $29.8$\,s & $3.89$\,s & $2.49$\,s & $0.79$\,s \\
 & $10^5$ & $27\,817$\,s & $3\,481$\,s & $1\,888$\,s & $389$\,s \\
\midrule
$T$ ($K_{-1}$) & $10^4$ & $31.9$\,s & $4.12$\,s & $3.36$\,s & $1.16$\,s \\
 & $10^5$ & $30\,050$\,s & $3\,762$\,s & $2\,287$\,s & $569$\,s \\
\midrule
$U(1)$ & $10^4$ & $53.5$\,s & $8.58$\,s & $61.3$\,s & $8.16$\,s \\
 & $10^5$ & $53\,348$\,s & $6\,905$\,s & -- & -- \\
\midrule
$H(1)$ & $10^4$ & $1.01$\,s & $0.15$\,s & $4.34$\,s & $0.93$\,s \\
 & $10^5$ & $304$\,s & $39.1$\,s & $1\,925$\,s & $345$\,s \\
\bottomrule
\end{tabular}
\vspace{\abovecaptionskip}
\caption{Timings of the reference methods in \texttt{holsum} with $8$ threads (CPU time summed over the threads, and wall time).}
\label{tab:appsref}
\end{table}

The theta splitting grows by a factor of $52$ to $69$ per decade, about $p^{1.75}$ to $p^{1.85}$, like the optimized
program for Khinchin's constant, and parallelizes as well (the ratio of CPU to wall time is $7.7$--$8.0$ at $10^5$ and
$10^6$ digits). The generic code is not far behind the optimized one: at $10^6$ digits, $\mathcal A$, whose states have
dimensions close to those for Khinchin's constant, takes $8\,435$ CPU seconds against $5\,416$. As expected from the
dimensions of the states, the cost grows from $\mathcal A$ (states of dimension $2$ and $3$) to $\mathcal C$ (odd
parts, $\Gamma(0,y)$ starting values) and to $T$ (harmonic numbers on both sides). For the vortex sum and for $H$ the
dual systems have order $4$ but few forcing constants, so that the dual terms are relatively cheap and the larger
split point $u_0 = \frac14$ is better; at $10^5$ digits, $u_0 = \frac1{16}$ is about $1.3$ times slower for the vortex
sum.

The zeta-value expansion grows by factors of $250$ to $1\,000$ per decade. At $10^5$ digits it is $8$ times slower than
the theta splitting for $\mathcal A$, which needs only even zeta values (Bernoulli numbers), and $30$ to $80$ times
slower for $\mathcal C$, $T$ and the vortex sum, where the odd zeta values dominate (for the vortex sum they are
needed to about $1.25p$ bits, since the coefficients grow like $\beta^j$). The exception is $H(1)$: its expansion needs
only even zeta values and, as its coefficients decay factorially, only about $p/\log p$ of them. It is faster than the
theta splitting at $10^5$ digits ($304$ against $419$ CPU seconds), but it grows by a factor of $300$ per decade
against $68$, so the crossover should be near $1.6\cdot10^5$ digits; it remains quasi-quadratic, since the Bernoulli
numbers involved have total size of order $p^2/\log p$. Euler--Maclaurin summation is the fastest method for the
logarithmic sums at $10^4$ digits, but at $10^5$ digits it needs $2$ to $12$ times more CPU time than the theta
splitting and $4.5$ to $20$ times more wall time, since its Bernoulli corrections are sequential. For the vortex sum its
direct part requires a ${}_2F_1$ value per term, and we did not run it beyond $10^4$ digits.

The values are
\begin{align*}
  \mathcal A &= 1.02313872642793929553508807697521809749214527936608\ldots, \\
  \mathcal C &= 0.78853056591150896106027632345455466647274966822328\ldots, \\
  e^{\mathcal C-1} &= 0.80939402054063913071793188059409131721595399242500\ldots, \\
  K_{-1} &= 1.74540566240734686349459630968366106729493661877798\ldots, \\
  U(1) &= 0.26882772504004215799420784633078761581126503327691\ldots, \\
  H(1) &= 1.47282823195618529629494738382314582532386592787930\ldots
\end{align*}
To our knowledge $\mathcal A$, $\mathcal C$ and $K_{-1}$ had previously been published to about $10^2$--$10^3$ digits
(OEIS b-files), except $K_{-1}$, which Bailey, Borwein and Crandall computed to $7\,350$ digits with the zeta-value
expansion of their Corollary~7 \cite{BBC97}. For other spacings of the vortex rings ($\lambda = 2/s$ with integer $s$,
so that $\beta = s^2$),
\begin{align*}
  U(2) &= 0.07770257288078689394748426153469590197848\ldots, \\
  U(\tfrac23) &= 0.48599721401691389437533080303387632179\ldots, \\
  U(\tfrac12) &= 0.71291908153312931408325016665401004450\ldots,
\end{align*}
which agree with direct summation of \eqref{eq:vortex} at $30$ digits.

The cross-checks are not a formality. A first version of the vortex program truncated the series for the starting
values too early (their $E$-part decays only like $(\beta u_0)^j/j!$), and the computations at different $u_0$ disagreed
beyond $75\,000$ digits; a first version of \texttt{holsum} lost about $\sqrt{\beta c}$ bits to cancellation in the
dual recurrences of the vortex sum, which the ball arithmetic revealed at $10^4$ digits.

\subsection{Further candidates}
The method applies to every constant in the tables of Flajolet and Vardi \cite{FV96} and of Finch \cite{Fin03}
that has a zeta-value expansion whose generating function $f$ is holonomic and satisfies (H) (with the terms $k < k_0$, $k_0 > 1/\rho$,
added separately): the H\"older means $K_{-n}$; constants of the form
$\sum_{k=k_0}^{\infty} F(1/k)$ with $F$ a product of logarithms, polylogarithms or inverse trigonometric functions (such as
$\sum_{k=1}^{\infty} \frac1k\arctan\frac1k$); and, through Theorem~\ref{thm:log}, sums $\sum_{k=k_0}^{\infty} (\log k)^j R(1/k)$ with $R$ rational.
Products $\prod_{k=k_0}^{\infty} F(1/k)$ reduce to such sums by taking logarithms when $\log F$ is holonomic.
Not every member of the class is new: sums of logarithms or inverse trigonometric or hyperbolic functions of a
rational function of $k$ are logarithms or arguments of products of gamma values at algebraic points, for instance
(A352619, A393668)
\[
  \sum_{k=1}^{\infty}\Bigl(\frac1k - \arctan\frac1k\Bigr) = \gamma + \arg\Gamma(1+i), \qquad
  \sum_{n=2}^{\infty}\arctanh\frac1{n^2} = \frac12\log\frac{\sinh\pi}{\pi},
\]
and these can be computed in quasilinear time.
The extra factor $\frac1k$ in $\sum_{k=1}^{\infty}\frac1k\arctan\frac1k$ destroys this product structure.

\subsubsection{Non-examples} Some classical sums with summands analytic at infinity are not in the class because the summand
is not holonomic, although their zeta-value expansions converge geometrically. Examples are the logarithm of the
Kepler--Bouwkamp constant and two constants from the OEIS (A342680, A336405),
\[
  \sum_{k=3}^{\infty}\log\cos\frac\pi k, \qquad \sum_{n=1}^{\infty}\Bigl(\tan\frac1n - \sin\frac1n\Bigr), \qquad
  \sum_{n=1}^{\infty}\log\Bigl(n\sin\frac1n\Bigr)
\]
(the derivative of $\log\cos$ is $-\tan$, which has infinitely many poles), sums of Clausen
functions $\operatorname{Cl}_2(2\pi/n)$, and sums involving $(1 - \frac2{n+1})^{2n}$, such as the oscillator strengths of the
hydrogen atom. Whether such sums admit subquadratic algorithms is open.

Sums of this kind also occur in physics, where they are usually truncated and the tail estimated by its leading
term. Besides the vortex-ring sum above (found by Vasilev and by Levy and Forsdyke, and summed with $500$ terms in
\cite{MS22}), examples are Matsubara-frequency sums of one-loop integrals at finite temperature \cite{HLS24}, whose
summands are arctangents of algebraic functions of the frequency index, and sums over Kaluza--Klein towers
\cite{Pet02}, whose summands involve $\arcsin^2$ of algebraic functions of the mode number. For small spacing $\lambda$,
the vortex-ring sum is a sum $\sum_{n=1}^{\infty} F(n\lambda)$ sampled finely, where the constants of Theorem~\ref{thm:general} grow.

%======================================================================
\section{Concluding remarks}\label{sec:concl}

The exponent $\frac32$ comes from the uncertainty principle of the Gaussian: $k_{\max}n_{\max} \approx qp\log2/\pi$ for
every split point, and every term costs at least $\Mul(p)$. Lowering it would require evaluating many terms at once.
On the dual side all points $c = \pi^2n^2/q^2$ share the irrational number $\pi^2$, and the bit-burst paths could share
their rational approximations, but the Taylor expansions differ with $n$. On the $k$-side, the series for different $k$
are values of one holonomic function of $1/k^2$ and $k^2u_0$; fast multipoint evaluation of such families might help.

In practice, the constant factor is decided by the order of the dual system, and identities that make the
coefficients hypergeometric (as Lemma~\ref{lem:dilog} does for Khinchin's constant) are worth looking for.

\begin{thebibliography}{99}
\bibitem{AG77} K.~Alladi and C.~Grinstead, On the decomposition of $n!$ into prime powers, J.\ Number Theory 9 (1977), 452--458.
\bibitem{AV84} T.~M.~Apostol and T.~H.~Vu, Dirichlet series related to the Riemann zeta function, J.\ Number Theory 19 (1984), 85--102.
\bibitem{BC21} K.~Belabas and H.~Cohen, \emph{Numerical Algorithms for Number Theory: Using Pari/GP}, Mathematical Surveys and Monographs 254, American Mathematical Society, 2021.
\bibitem{BBC97} D.~H.~Bailey, J.~M.~Borwein and R.~E.~Crandall, On the Khintchine constant, Math.\ Comp.\ 66 (1997), 417--431.
\bibitem{BGMWZ13} J.~M.~Borwein, M.~L.~Glasser, R.~C.~McPhedran, J.~G.~Wan and I.~J.~Zucker, \emph{Lattice Sums Then and Now}, Cambridge University Press, 2013.
\bibitem{BGS07} A.~Bostan, P.~Gaudry and \'E.~Schost, Linear recurrences with polynomial coefficients and application to integer factorization and Cartier--Manin operator, SIAM J.\ Comput.\ 36 (2007), 1777--1806.
\bibitem{BB87} J.~M.~Borwein and P.~B.~Borwein, \emph{Pi and the AGM}, Wiley, 1987.
\bibitem{Boy19} K.~N.~Boyadzhiev, A special constant and series with zeta values and harmonic numbers, arXiv:1903.11141, 2019.
\bibitem{BZ10} R.~P.~Brent and P.~Zimmermann, \emph{Modern Computer Arithmetic}, Cambridge University Press, 2010.
\bibitem{BBC00} J.~M.~Borwein, D.~M.~Bradley and R.~E.~Crandall, Computational strategies for the Riemann zeta function, J.\ Comput.\ Appl.\ Math.\ 121 (2000), 247--296.
\bibitem{BM80} R.~P.~Brent and E.~M.~McMillan, Some new algorithms for high-precision computation of Euler's constant, Math.\ Comp.\ 34 (1980), 305--312.
\bibitem{CC20} B.~Candelpergher and M.-A.~Coppo, Laurent expansion of harmonic zeta functions, preprint, HAL hal-02444740.
\bibitem{CN62} K.~Chandrasekharan and R.~Narasimhan, Functional equations with multiple gamma factors and the average order of arithmetical functions, Ann.\ of Math.\ 76 (1962), 93--136.
\bibitem{CN81} P.~Cassou-Nogu\`es, Applications arithm\'etiques de l'\'etude des valeurs aux entiers n\'egatifs des s\'eries de Dirichlet associ\'ees \`a un polyn\^ome, Ann.\ Inst.\ Fourier 31 (1981), no.~4.
\bibitem{CHTZZ07} H.~Cheng, G.~Hanrot, E.~Thom\'e, E.~Zima and P.~Zimmermann, Time- and space-efficient evaluation of some hypergeometric constants, Proc.\ ISSAC 2007, 85--91.
\bibitem{CS67} S.~Chowla and A.~Selberg, On Epstein's zeta-function, J.\ Reine Angew.\ Math.\ 227 (1967), 86--110.
\bibitem{CC90} D.~V.~Chudnovsky and G.~V.~Chudnovsky, Computer algebra in the service of mathematical physics and number theory, in \emph{Computers in Mathematics}, Lecture Notes in Pure and Appl.\ Math.\ 125, Dekker, 1990, 109--232.
\bibitem{Cof14} M.~W.~Coffey, Series representations for the Stieltjes constants, Rocky Mountain J.\ Math.\ 44 (2014),
443--477.
\bibitem{Cof12} M.~W.~Coffey, Series representation of the Riemann zeta function and other results: complements to a paper of Crandall, Math.\ Comp.\ 83 (2014), 1383--1395.
\bibitem{CRVZ00} H.~Cohen, F.~Rodriguez Villegas and D.~Zagier, Convergence acceleration of alternating series, Experiment.\ Math.\ 9 (2000), 3--12.
\bibitem{Cra12} R.~E.~Crandall, Unified algorithms for polylogarithm, $L$-series, and zeta variants, preprint, 2012.
\bibitem{Dok04} T.~Dokchitser, Computing special values of motivic $L$-functions, Experiment.\ Math.\ 13 (2004), 137--149.
\bibitem{EKR11} W.~Ebeling, W.-D.~Kraeft and G.~R\"opke, Bound states in Coulomb systems -- old problems and new solutions, arXiv:1110.1962, 2011.
\bibitem{Eli95} E.~Elizalde, \emph{Ten Physical Applications of Spectral Zeta Functions}, Lecture Notes in Physics m35, Springer, 1995.
\bibitem{Ewa21} P.~P.~Ewald, Die Berechnung optischer und elektrostatischer Gitterpotentiale, Ann.\ Phys.\ 369 (1921), 253--287.
\bibitem{Fin03} S.~R.~Finch, \emph{Mathematical Constants}, Cambridge University Press, 2003.
\bibitem{FGS05} P.~Flajolet, S.~Gerhold and B.~Salvy, On the non-holonomic character of logarithms, powers, and the $n$th prime function, Electron.\ J.\ Combin.\ 11(2) (2004--2005), \#A2.
\bibitem{Fin18} S.~R.~Finch, \emph{Mathematical Constants II}, Cambridge University Press, 2018.
\bibitem{FGD95} P.~Flajolet, X.~Gourdon and P.~Dumas, Mellin transforms and asymptotics: harmonic sums, Theoret.\ Comput.\ Sci.\ 144 (1995), 3--58.
\bibitem{FV96} P.~Flajolet and I.~Vardi, Zeta function expansions of classical constants, manuscript, 1996.
\bibitem{FLINT} The FLINT team, FLINT: Fast Library for Number Theory, \url{https://flintlib.org}.
\bibitem{Gha26} R.~K.~Ghazi, khinchin-fast: Khinchin's constant to $2\,000\,000$ certified decimal digits, Zenodo, 2026, \url{https://doi.org/10.5281/zenodo.22134905}.
\bibitem{GLS03} R.~Garun\v{k}\v{s}tis, A.~Laurin\v{c}ikas and J.~Steuding, An approximate functional equation for the Lerch zeta function, Math.\ Notes 74 (2003), 469--476.
\bibitem{Gau05} W.~Gautschi, The Hardy--Littlewood function: an exercise in slowly convergent series, J.\ Comput.\ Appl.\ Math.\ (2005), 249--254.
\bibitem{GM85} W.~Gautschi and G.~V.~Milovanovi\'c, Gaussian quadrature involving Einstein and Fermi functions with an application to summation of series, Math.\ Comp.\ 44 (1985), 177--190.
\bibitem{HLS24} D.~Harnett, S.~Li and T.~G.~Steele, Numerically computing finite temperature loop integrals using pySecDec, Eur.\ Phys.\ J.\ A (2024), arXiv:2401.05258.
\bibitem{HP98} B.~Haible and T.~Papanikolaou, Fast multiprecision evaluation of series of rational numbers, Proc.\ ANTS-III, LNCS 1423, Springer, 1998, 338--350.
\bibitem{HvdH21} D.~Harvey and J.~van der Hoeven, Integer multiplication in time $O(n\log n)$, Ann.\ of Math.\ 193 (2021), 563--617.
\bibitem{vdH99} J.~van der Hoeven, Fast evaluation of holonomic functions, Theoret.\ Comput.\ Sci.\ 210 (1999), 199--215.
\bibitem{Inc26} E.~L.~Ince, \emph{Ordinary Differential Equations}, Longmans, Green and Co., 1926.
\bibitem{IK02} M.~Iosifescu and C.~Kraaikamp, \emph{Metrical Theory of Continued Fractions}, Kluwer, 2002.
\bibitem{Joh17} F.~Johansson, Arb: efficient arbitrary-precision midpoint-radius interval arithmetic, IEEE Trans.\ Comput.\ 66 (2017), 1281--1292.
\bibitem{Joh14} F.~Johansson, Evaluating parametric holonomic sequences using rectangular splitting, Proc.\ ISSAC 2014, 256--263; arXiv:1310.3741.
\bibitem{Joh21} F.~Johansson, Rapid computation of special values of Dirichlet $L$-functions, arXiv:2110.10583, 2021.
\bibitem{Kar95} E.~A.~Karatsuba, Fast calculation of the Riemann zeta function $\zeta(s)$ for integer values of the argument $s$, Problems Inform.\ Transmission 31 (1995), 353--362.
\bibitem{Kar98} E.~A.~Karatsuba, Fast evaluation of the Hurwitz zeta function and Dirichlet $L$-series, Problems Inform.\ Transmission 34 (1998), 342--353.
\bibitem{Kal88} S.~Kalpazidou, Khintchine's constant for L\"uroth representation, J.\ Number Theory 29 (1988), 196--205.
\bibitem{Khi35} A.~Khintchine, Metrische Kettenbruchprobleme, Compositio Math.\ 1 (1935), 361--382.
\bibitem{Kir02} K.~Kirsten, \emph{Spectral Functions in Mathematics and Physics}, Chapman \& Hall/CRC, 2002.
\bibitem{Lev36} P.~L\'evy, Sur le d\'eveloppement en fraction continue d'un nombre choisi au hasard, Compositio
Math.\ 3 (1936), 286--303.
\bibitem{Loc64} G.~Lochs, Vergleich der Genauigkeit von Dezimalbruch und Kettenbruch, Abh.\ Math.\ Sem.\ Univ.\
Hamburg 27 (1964), 142--144.
\bibitem{Lav68} A.~F.~Lavrik, Approximate functional equations of Dirichlet functions, Izv.\ Akad.\ Nauk SSSR Ser.\ Mat.\ 32 (1968), 134--185.
\bibitem{MS22} E.~Masroor and M.~A.~Stremler, On the topology of the atmosphere advected by a periodic array of axisymmetric thin-cored vortex rings, Regul.\ Chaotic Dyn.\ 27 (2022), 183--197.
\bibitem{Mah28} K.~Mahler, \"Uber einen Satz von Mellin, Math.\ Ann.\ 100 (1928).
\bibitem{Mon10} H.~Monien, Gaussian quadrature for sums: a rapidly convergent summation scheme, Math.\ Comp.\ 79 (2010), 857--869.
\bibitem{Mez10} M.~Mezzarobba, NumGfun: a package for numerical and analytic computation with D-finite functions, Proc.\ ISSAC 2010, 139--146.
\bibitem{Mez16} M.~Mezzarobba, Rigorous multiple-precision evaluation of D-finite functions in SageMath, arXiv:1607.01967, 2016.
\bibitem{OEIS} OEIS Foundation Inc., The On-Line Encyclopedia of Integer Sequences, \url{https://oeis.org}.
\bibitem{OAI26} OpenAI, Integer multiplication below $n\log n$, OpenAI Math Release preprint, September 2026,
\url{https://github.com/openai/math/tree/main/preprints/Integer-multiplication-below-n-log-n-September-23-2026}.
\bibitem{Pet98} M.~Peter, Dirichlet series associated with polynomials, Acta Arith.\ 84 (1998), 245--278.
\bibitem{Pet02} F.~J.~Petriello, Kaluza--Klein effects on Higgs physics in universal extra dimensions, arXiv:hep-ph/0204067, 2002.
\bibitem{Rie59} B.~Riemann, Ueber die Anzahl der Primzahlen unter einer gegebenen Gr\"osse, Monatsber.\ Berliner Akad., 1859.
\bibitem{Rub05} M.~Rubinstein, Computational methods and experiments in analytic number theory, in \emph{Recent Perspectives in Random Matrix Theory and Number Theory}, LMS Lecture Note Ser.\ 322, Cambridge University Press, 2005, 425--506.
\bibitem{Sta80} R.~P.~Stanley, Differentiably finite power series, European J.\ Combin.\ 1 (1980), 175--188.
\bibitem{Sar84} P.~Sargos, Prolongement m\'eromorphe des s\'eries de Dirichlet associ\'ees \`a des fractions rationnelles de plusieurs variables, Ann.\ Inst.\ Fourier 34 (1984), no.~3, 83--123.
\bibitem{Sch51} J.~Schwinger, On gauge invariance and vacuum polarization, Phys.\ Rev.\ 82 (1951), 664--679.
\bibitem{SW59} D.~Shanks and J.~W.~Wrench, Khintchine's constant, Amer.\ Math.\ Monthly 66 (1959), 276--279.
\bibitem{MWKdigits} E.~W.~Weisstein, Khinchin's constant digits, MathWorld, \url{https://mathworld.wolfram.com/KhinchinsConstantDigits.html}, accessed November 2026.
\end{thebibliography}

\end{document}
